\chapter{History of the Sri Lankan Malays}

% \begin{motto}
%    [T]he linguist who makes theories about language influence but neglects to account for the socio-cultural setting of the language contact leaves his study suspended, as it were, in mid-air. \citep[4]{Weinreich1953}
% \end{motto}
 
Sri Lanka Malay is a contact language. The socio-cultural setting is very important for analyses of language contact \citep{Weinreich1953, ThomasonEtAl1988}. This chapter shall present the
different periods of contact between Malays and Lankans and
analyze the socio-historical setting in each period.

The first significant number of Malays arrived in
Ceylon in the 17$^{th}$ century. However, the contact between the
islands of Ceylon and the South East Asian Archipelago (SEAA)
predates this by many centuries. We can divide the history of
contact between the Malays and the Lankans drawing on the
following major historical events:

\begin{itemize}
    \item The first Malays brought to Ceylon by the Dutch (\Tilde 1640)
    \item The ousting of the Dutch by the British and the creation of a Malay regiment (\Tilde 1796)
    \item The disbandment of said regiment (1873)
    \item independence of Ceylon from the United Kingdom (1948)
\end{itemize}

Five periods can thus be distinguished:
a pre-Dutch period (-1640, Section \ref{sec:slmbg:HistorybeforethearrivaloftheDutch} p.\pageref{sec:slmbg:HistorybeforethearrivaloftheDutch}),
the Dutch period (1640-1796, Section \ref{sec:slmbg:HistoryduringDutchrule} p.\pageref{sec:slmbg:HistoryduringDutchrule}),
the period of the Regiment under the British (1796-1873, Section \ref{sec:slmbg:HistoryduringBritishrule} p.\pageref{sec:slmbg:HistoryduringBritishrule})\footnote{Technically speaking, the regiment was founded some years later than 1796.},
the post-regimental history and the post-colonial history (1873-today, Section \ref{sec:slmbg:PostregimentalHistory} p.\pageref{sec:slmbg:PostregimentalHistory}).

I will briefly sketch major events during each period before discussing demographic and sociological developments that accompanied them. Domains which encompass more than one period are
the cultural history (Section \ref{sec:slmbg:CulturalHistory} p.\pageref{sec:slmbg:CulturalHistory}),
the Malay settlements in Sri Lanka \ref{sec:slmbg:Malaysettlements} (p.\pageref{sec:slmbg:Malaysettlements}),
and the relation to the other Islamic group, the `Moors' (Section \ref{sec:slmbg:RelationtotheMoors}  p.\pageref{sec:slmbg:RelationtotheMoors}),
The development of the language will be discussed in Section \ref{sec:slmbg:LanguageIssues}, p.\pageref{sec:slmbg:LanguageIssues}.

\section{Before the arrival of the Dutch}\label{sec:slmbg:HistorybeforethearrivaloftheDutch} 

\subsection{Political history}\label{sec:slmbg:preDutch:Political history} 
The first known people to settle in Ceylon were the Veddah, who
live today only in remote pockets of jungle in the central
highlands.
It can surely be affirmed that the Veddah have not had any
linguistic relation on the present day SLM because they were
already pushed back into the jungles by the time the Malays entered
Lanka.

The Veddah could not withstand the arrival of the populations from the Indian mainland about 500 BC \citep[1]{GairEtAl1997} that announces a second phase of settlement, resulting in the present cultures of Sinhalese and Ceylon Tamils.\footnote{While for the Indo-Aryans, who were to become the Sinhalese later, 500 BC seems a reasonable estimation, for the Dravidians ``there is not firm evidence as to when [they] first came to the island, but come they did from very early times.'' \citep[12]{Desilva1981}.}

From the 7^{th} century AD onwards, Arab traders settled down in  India \citep[51]{Codrington1926}. The Arab traders in India became acquainted with local tongues. Some of their Tamil-speaking descendants moved to Ceylon from the seventh \citep[32]{Nuhman2007}, eighth \citep[21]{Hussein2007} or tenth \citep[72]{Desilva1981} century onwards. They are the ancestors of most of today's Tamil speaking Muslims in Lanka.

% 
% The  relations Ceylon and South East Asia (SEA) in the precolonial period are not clear. People from the South East Asian Archipelago (SEAA) settled in Madagascar from the 2$^{nd}$ century AD and it is reasonable to suppose that they stopped over in Ceylon on their way. However, they left no traces. But when they went back to the SEAA, they took Buddhism with them, and there are iconographical finds in Western Java, dated to the 5$^{th}$ or 6$^{th}$ century AD \citep{Dupont1959}  and a 6$^{th}$ century inscription found in Central Java \citep{DeCasparis1961}  that point to a more than superficial contact between the island peoples of Ceylon and the SEAA \citep[5]{Bichsel}.
% 
% According to  \citet[11]{Jayasuriya2002},
% 
% \begin{quotation}  
%     Contact between Malaysia, Indonesia and Sri Lanka dates back at least to the 8$^{th}$ century when the Sailendra dynasty which originated from Java, ruled over Malaya, Sumatra, Java, Bali and Borneo. Little is known about the dynasty but the Buddhist monuments such as Borobudur were constructed during their rule. The   Sailendra royalty were Mahayana Buddhists whose religious texts were written in Sanskrit \el. There were also genealogical connections between the Sailendra dynasty and the Sinhalese royalty who ruled Sri Lanka.  \citep[11]{Jayasuriya2002}
% \end{quotation}
% 
% What is surprising in the analysis of Jayasuriya is that the Sailendra were Mahayana Buddhists. This does not fit to well with a close link with Ceylon, the religious center of Theravada Buddhism. If the religion of the rulers of Java is to provide evidence, then their following of Mahayana has to be regarded as countering claims of a link. Genealogy on the other hand is a good indicator of close contact, unfortunately Jayasuriya does not provide  the exact relation between the Sailendra and the Sinhalese royalty.


While we have evidence of cultural Lankan influence in the SEAA from the
first millennium AD \citep{Dupont1959,DeCasparis1961}, the earliest source of Malay culture in Sri
Ceylon dates from a later period. It is found in old Sinhala
literature from the periods of Polonnaruwa and the Dambadeniya
(together 1098-1293) which contains words of Malay etymology. More
recent literature from the 15$^{th}$ century states that the SEAA was
renowned for spices, perfumes and other commodities
\citep[11]{Jayasuriya2002}. While this testifies  to some sort of
contact between Ceylon and the SEAA, the impact seems to have been
rather minor and restricted to some loanwords \citep{Gunasekara1891,Sannasgala1976}. Cultural
elements that would point to a more intimate contact could be
outrigger canoes \citep{Tennent1859}.

The peaceful exchange of goods and religion found an end in
1247, when Malay forces from the peninsula made an attempt
to invade Sri Lanka \citep[67]{Desilva1981}. Chandrabh\=anu of T\=ambralinga attacked the Polonnaruwa kingdom in Central Sri Lanka. He was defeated and withdrew to the Jaffna peninsula, where he succeeded in securing the throne for himself. Later, he tried a second invasion of Polonnaruwa, which failed and cost him his life, ending the short Malay intermezzo on the island. Today's
place names \em Caavakacceeri \em (Javaka-settlement) and Caavakotte (Javaka
Fort) witness this invasion.\footnote{\em Java \em is taken as a
generic term for the SEAA in Ceylon \citep[37]{Hussainmiya1987}.} Even if this invasion
had a period of Malay rulership as a consequence, there is no link to the present day Malay communities \citep[cf.][685]{AdelaarEtAl1996}. An invasion that should have a much bigger impact on Lankan society was threatening from another continent: the Portuguese.

% \begin{quote}
% 	``Less obvious than the Indian influence, but over the centuries just as important, was the influence form South-East Asia resulting from Sri Lanka's strategic location athwart the main sea-routes of the Indian Ocean.''\citet[4]{Desilva1981}
% \end{quote}

 
Driven by the desire for the lucrative spice trade, the Portuguese sent out expeditions to secure a base in the SEAA. In 1511 the Portuguese gained hold of Malacca on the Malaysian peninsula. By that time, Islam had made considerable inroads in SEAA and largely replaced Buddhism as the predominant religion. On their way from the Cape to the SEAA, Ceylon was a good base for the Portuguese, and in 1518  they established a fort in Colombo. Within the following decades, they gained control of the coast but could not penetrate the central highlands, where the Kandyan king would withstand European conquest until the 19$^{th}$ century.

A major interest of the Portuguese was to spread Christian faith.
They kept other religions in very low regard, especially Islam, and brought about the expulsion of Muslims from the K\=o\dott\dott e kingdom \citep[102]{Desilva1981}.
\citet[34f]{Hussainmiya1990} therefore questions whether the Portuguese would have
brought (predominantly Muslim) soldiers from the SEAA to fight for
them. \citet[6f.]{Bichsel} agrees that there is no historical
evidence for a sizeable Malay community in Sri Lanka and doubts
that the Portuguese would have tolerated the spread of Malays in
Sri Lanka.  However, \citet[26]{Powell1973} mentions that some \em
Javas \em  and \em Kaffirs \em (Africans) fought for the
Sinhalese king Raju in 1587 and states that they had deserted from
the Portuguese, indicating that some Malays were already found in the Portuguese forces.

% Beside the historical sources, some scholars have also tried to find linguistic evidence for contact between the Malays and the Portuguese.  \citet[5f]{Jayasuriya2002} affirms that there is evidence for contact between Sri Lanka Portuguese Creole (SLPC, \citet{Smith1979}) and Malayan Dutch. She seems to imply that this must have taken place during the Portuguese period, but a later contact between SLPC speakers and Malay Dutch speakers is also conceivable. The SLPC speakers stayed on the island after the Portuguese defeat and could easily have picked up Dutch Malayan words during Dutch rule.

\subsection{Demographic history}\label{sec:slmbg:DemographichistorybeforearrivaloftheDutch}
The study of Malay demography before the arrival of the Dutch is hampered by the paucity of subjects to study. There
was no important settlement of Malays in Ceylon before the arrival
of the Dutch. Chandrabh\=anu's invasion did not last long enough
to give birth to a Malay community in Ceylon.  If any Malays were
left behind after the invasion, they would have merged in their
coreligionist community of Sinhalese Buddhists. The Malays brought
by the Portuguese do not seem to have been very numerous, and we
do not have any information whatsoever about  where they had come from or
what languages they might have spoken. These Malays are, however,
the first to desert to the Kandyan King and may have had a certain
influence on the Malay community in Kandy.


\section{Dutch rule}\label{sec:slmbg:HistoryduringDutchrule}
The Dutch had gained considerable power and prosperity in the 17$^{th}$ century and projected their power to the Indian Ocean. By 1658, they had managed to oust the Portuguese from Ceylon with the help from the Kandyan King. The latter, however, had to find out that his ``allies'' had come to stay and that he had bartered Portuguese occupation of the coastline  for Dutch occupation of the coastline \citep[133]{Desilva1981}.

\subsection{Political history}\label{sec:slmbg:Dutch:Political history}
Like the Portuguese before them, the Dutch had taken hold in the SEAA before trying to get hold of Ceylon.  In their possession in the SEAA, they had met peoples with excellent soldiers, who the Dutch recruited into their service. These troops were already fighting on the Dutch side during the war with the Portuguese  \citep[44]{Hussainmiya1990}.
Later, they formed an important part of the Dutch military presence in Ceylon.  There were garrisons in the main coastal towns of Colombo, Galle, Chilaw and Trincomalee.

The Eastern troops had several advantages over European troops: First, they were used to tropical climate and fights in the jungle \citep[46]{Hussainmiya1990} and,  second, they were not likely to side with the Sinhalese or Tamils in any future rebellion  given the cultural and religious differences. \citet[55]{Hussainmiya1987} reports that the Dutch referred to these troops as \trs{Oosterlingen}{Easterlings} or \trs{Javaans}{Javanese},  which is a much better name than `Malay', which the British authorities would  later choose as a generic term. Actually, the majority of the immigrants were \em not \em Malay, and were definitely not from the Malaysian peninsula. In the Ceylonese context, however, the word `Malay' is reasonably unambiguous (it can only refer to the group under discussion), and it is common scientific practice to refer to this immigrant group as `Malay' for want of a better word, which is also what I will do in the remainder of this book.

A regular supply of new recruits from Batavia was provided each year, as well as reinforcements in urgent cases \citep[48]{Hussainmiya1987}. The importance of the Malays for the Dutch army can be seen from the fact that ``[a]bout two thirds of the fighting troops were Malays'' in 1764 \citep[37]{Powell1973}, a year in which  the absolute number of Malay soldiers  also reached the record of 2500 \citep[49]{Hussainmiya1990}. This coincided with the creation of the Free Javanese Company \citep[52]{Hussainmiya1990}. From 1785 onwards, local, i.e. Ceylonese, civilian Malays were recruited into the Dutch army to replace the disbanded regiment of Free Moors  \citep[54]{Hussainmiya1987}.

Beside the soldiers, who represented the numerically  most important group, some other groups   from the SEAA were also brought to Ceylon: exiles from rebellious provinces, convicts, and slaves.

In order to stay in control of the islands of the SEAA, the Dutch
deported rebellious local regents to South Africa and Ceylon.   Ceylon
was preferred because it was closer to the SEAA; very
unaccommodating princes would be exiled farther to Africa. The
Dutch having first taken possession of the former Portuguese
colonies in the Moluccas and the lesser Sunda islands, the first
exiles came from there. In the late 17$^{th}$ century, a war of
succession in Java had the Dutch deport some Javanese princes to
Ceylon \citep[38]{Hussainmiya1990}.

Convicts from the SEAA were brought to Ceylon, ``serving sentences according to the gravity of their crime ranging from being kept in chains to the permission to engage in handicraft''  \citep[44]{Hussainmiya1990}.
In 1782, the Dutch allowed minor offenders to join the ranks of the military. Other convicts were set free and were allowed to settle in Ceylon.   The  former convicts  thus formed part of the early Malay population in Ceylon and the `Malay' community of Sri Lanka can be said to owe its origins in part to these Malay convict settlers.

A third group were slaves and servants brought mainly from the Moluccas and Lesser Sunda islands to serve the aristocratic exiles and the Dutch. Like the convicts, some of them joined the army \citep[48]{Hussainmiya1990}. They were not a very numerous group and did not constitute a distinguishable separate community.

Particularly after the middle of the 18$^{th}$ century, the proportion of free Malays grew. From 1769 onwards, the free Malays had  their own category for tax collecting purposes \citep[9]{Bichsel}, which indicates that their number could not be neglected. This comprised discharged soldiers and descendants of political exiles. Their occupations included gardening, rattan weaving and petty trade. Some  were also employed as domestic servants. These Malays counted as foreigners in Ceylon and had to perform \em Uliyam \em (soccage),\footnote{For more on soccage,  see \citet[31]{Ludowyk1966}, \citet[149f]{Codrington1926}, \citet[164]{Desilva1981}.} a duty which the  political exiles and soldiers were exempt from \citep[48f]{Hussainmiya1990}.

% \begin{quote}
% 	``The Dutch Government, which in October [1795] had resolved to resist to the last, had at its disposal 1617 men, composed of 845 Europeans and 772 Malays.'' \citep[157]{Codrington1926}
% \end{quote}


\subsection{Demographic history}\label{sec:slmbg:DemographichistoryunderDutchreign}
The origin of the  four different groups of immigrants (Exiles, convicts, slaves, and soldiers)  will be discussed in this section.

The first exiles came from the Moluccas and the Lesser Sunda Islands.
Later, there were also exiled Javanese princes \citep[38f]{Hussainmiya1990}. Hussainmiya further cites Bacan, Tidore, Timor, Madura and Sumatra. The Javanese
succession wars in the late 17$^{th}$ century brought more Javanese
exiles to Ceylon.  Most of the exiles had their families with them
at deportation time or they joined them later.  \citet[39]{Hussainmiya1990} gives 200 as a likely number for the exiles at the end of the 18$^{th}$ century, about 10\% of the Malay population at the time. Most of the exiles lived in the 4 main coastal towns Colombo, Galle, Trincomalee and Jaffna.

While we have a certain knowledge of how many exiles came from which region, the case is more difficult for the convicts and the slaves \citep[47]{Hussainmiya1987}. The number of slaves is uncertain but small, and they came from the Moluccas and the Lesser Sunda islands \citep[48]{Hussainmiya1990}; the provenance of the convicts is unknown. In 1751 there were 131 convicts serving their sentence; the number of convicts having joined the army or regained freedom is unclear.


\begin{sidewaysfigure}
 \centering
\includegraphics[width=\textwidth]{pics/Dutchimmigration} 
 \caption[Provenance of Malay immigrants during Dutch rule]{The provenance of the Malay immigrants during Dutch rule. The size of the arrows corresponds to the relative strengths of the contingents. The routes represented on the graphics are not necessarily those employed by the immigrants, with the exception of the stop-over of soldiers   in Batavia/Jakarta.  Small arrows represent exiles, whose precise numbers are not known.}
\label{fig:ImmigrationDutch}
\end{sidewaysfigure}

 
The most important contingent of Malays brought by the Dutch were, however, the  soldiers.  They were mainly recruited in Batavia (Jakarta). At that time, Batavia was a port city with many different peoples mingling.  There are records of the following ethnic groups serving in the military: Ambonese, Bandanese, Balinese, Buginese, Javanese, Madurese, Malays \citep{Bichsel,Raben2000}.  In Ceylon, they were referred to as \em  Javaans\em, but we see now that this is dramatically simplified given the ethnic diversity of the newcomers. These ``Neo-Indonesians'', as \citet[57]{Hussainmiya1987} calls them, formed the base for the Malay society in Ceylon.  Already in Batavia, they are likely to have developed a group identity while waiting to be sent to the battlefield  \citep[174]{Vlekke1943}. This identity and cohesion would have been further strengthened  in the unknown foreign country.

In Batavia, there were different \em kampungs \em for the different ethnic groups \citep[451ff]{DeHaan1935}. %\citep[174]{Vlekke1943}.
But this does not mean that there was no contact between the different groups. The \em kampungs \em were not autarkic,\footnote{\citet[97]{Raben2000} criticizes the uncritical interpretation of Dutch rules regarding segregation by \citet{DeHaan1935} and argues that ``a yawning gulf existed between law and implementation.'' Also see \citet[31]{Abeyasekere1987}.} and one had to interact with other communities and make oneself understood. To this end, inhabitants employed a language of interethnic communication, a Malay-based lingua franca in their case (\citet[470]{DeHaan1935}, \citet[102]{Raben2000}). All over the SEAA, the use of some lingua franca is common, which varies from region to region \citep{AdelaarEtAl1996}. The nature of this language will be discussed in more detail in Section \ref{sec:slmbg:SourcesintheSEAA}.


Most of the nationalities present in Batavia were Muslims, but not all of them, as  \citet[58]{Hussainmiya1987} states:

\begin{quote}
    In Dutch times, not all the easterners who came to Sri Lanka were the followers of Islam. It is particularly difficult to establish the religious background of the Amboinese [sic], Balinese, and even Javanese, because among the first group, there  were a considerable number of Christians, while most of the Balinese belonged to Hindu Religion \el Nevertheless, it is quite clear that a large majority of the early Malays were followers of Islam. \citep[58]{Hussainmiya1987}
\end{quote}

The number of Malays serving varied in function of the military needs of the Dutch. Under normal conditions, there were 800 men in service, divided into 10 companies.  But sometimes the Dutch need for soldiers was so dramatic that the \em kampungs \em
were virtually depopulated \citep[56]{Hussainmiya1987}. In 1764 the number of Malay soldiers rose from 800 to 2500 in anticipation of a war with the Kandyan king. In 1795/6 there were probably around 1400 Malays serving, comprising recruited
civilians. For 1800, \citet[49f]{Hussainmiya1990} estimates the number of Malays in Ceylon  at around 2400, possibly their highest figure \citep[54]{Hussainmiya1987} because the Dutch were anticipating a British attack and had stocked up.
 
 
To sum up, during the Dutch period, there were two main groups. First there was the high prestige group of the exiles from the Moluccas and the Lesser Sunda Islands. This group made up about 10\% of the Malay population. Second, there were the soldiers, who came from diverse ethnic backgrounds, but who most likely communicated in some kind of vehicular Malay already developed back in Batavia. These soldiers formed the most important group, and they are likely to have absorbed the smaller groups of the slaves and convicts, who could not affirm their own identity given their low prestige and often joined the ranks of the soldiers where they had to adapt to the language already in use.



During Dutch rule, a shift in the composition of the Malay population in Ceylon took place. While there were many exiles from the Moluccas and Lesser Sundas at first, they were quickly outnumbered by exiles from the Javanese succession war and the Batavian soldiers. It seems that the Eastern community was absorbed by the Javanese community at the close of the 18$^{th}$ century.  Immigrants of the late 18$^{th}$ century would find an established society of ``Ceylon Malays'' with strong roots in Batavia, into which they would quickly integrate \citep[55-57]{Hussainmiya1987}.


\subsection{Sociological history}\label{sec:slmbg:SociologicalhistoryunderDutchrule}

\citet{Hussainmiya1987,Hussainmiya1990} provides some sociological information for the Malays under Dutch rule, but at times this information is contradictory. These contradictions concern the degree of isolation of the exiles, the extent of contact with the local population, and the possibility to practice their faith.

The two main groups of the exiles and the soldiers formed the bulk of the Malay population. Slaves and convicts were much less important in number and had regular connection with the soldiers, whose ranks they could join to gain their freedom \citep[48]{Hussainmiya1990}. It appears that the exiles formed  an elite of spiritual leadership, while the soldiers and the associated groups contributing  most of the souls.

The exiles were not allowed to communicate with each  other, but in how far this was enforced remains questionable given that marriages between children of exiles were contracted and attendance at celebrations such as births was also high \citep[43f]{Hussainmiya1990}. Some exiles were given ranks in the military in order to raise the moral of the Malay troops \citep[50]{Hussainmiya1990}, linking the exiles with the soldier population, and probably also communicating with each other.

% \citet[41]{Hussainmiya1990} affirms that the exiles were rather poor and had to sell jewels and other precious things in order to meet their expenses. Others were given land, but this could hardly earn them a living \citep[45]{Hussainmiya1990}.
% However, the situation of an exile should still have been a lot better than a peasant's situation of the period who did not have jewellery to sell and the alleged poverty of the exiles has to be seen in comparison to their previous standard of living.

While the exiles and the soldiers would see their stay in Ceylon as temporary, the free Malays had no plans to return to the SEAA. They are likely to have invested more in building up an independent existence in Ceylon.  Free Malays engaged in commercial activities such as gardening, rattan weaving and petty trade \citep[53]{Hussainmiya1987}. This  brought them into contact with the other Malay groups, but also with the other ethnic groups in Ceylon.
 
As far as the contact between the Malays and Lankan peoples is concerned, the situation varies. As for the exiles, normally, only the head of a family was considered such; the other members of the family could move about freely. Given the low number of exiles and the interdiction to communicate with other exiles, contact with locals seems probable. On the other hand, the exiles had no plans of staying in Ceylon for a long time, which makes them less likely to intermingle with the local population \citep[cf.][160]{Muysken2001sla}. Contact could also be inhibited by the Malay's residence inside the Dutch forts or in houses guarded by sentries \citep[43]{Hussainmiya1987}.
As for the soldiers, they fought for the Dutch against the local population. This makes friendly contact unlikely.  Different religious beliefs did not help the issue. \citet[cf.][51]{Hussainmiya1987} adduces that the Sinhalese even feared the Malays, and it seems reasonable to assume that the same should be true for the Hindu Tamils.
The Malays seem to have participated in the local cultural life, though; cockfighting was very popular, for instance \citep[131]{Pieris1918}.

The business within a Malay regiment could be carried out in Malay or in Dutch. Provisions and housing were also provided by the Army. Existing sources suggest that the need for interaction with the local population only arose in order to find a suitable wife.\footnote{However,
 official sources tend to exclude unofficial activities such as smuggling \citep[68]{Pieris1918}, gambling (\citet[131]{Pieris1918}, \citet[106]{Schweitzer1931}), and prostitution, which might have brought the groups into contact without this being reflected in official documents.}
Even this would not be necessary if the soldier had brought a wife from the SEAA, which was common \citep[cf.][50]{Hussainmiya1987}.


Sources on religious life are contradictory. On the one hand, the Dutch government in Ceylon  had banned public Islamic religious ceremonies, which was why they were often held in secrecy  \citep[46]{Hussainmiya1987}. On  the other hand, we have sources stating that from 1783 onwards, Malay mosques existed \citep[9]{Bichsel}. \citet[11]{Bichsel} affirms that the Malays were a clearly different community with distinct ethnic identity and own places of worship.

\subsection{Early Intermarriage}\label{sec:slmbg:EarlyIntermarriage}
It seems that already in the early ages of Malay presence in Ceylon, the Malays intermarried with the local population. Christopher Schweitzer, an employee of the VOC notes:

\begin{quote}
The wives, who in part are Ambonese, in part Sinhalese and Malabarian may not say anything [against the stripping of their ornaments]. \citep[106]{Schweitzer1931}, [my translation]\footnote{\el  die Weiber, welche theils Ambonesin, Singulesin und Malabarin, dörffen nichts darzusagen.}
\end{quote}

and also shortly after the British conquest, the commandant of Kalutara would write:

\begin{quote}
     As far as I can learn there is little damage to be apprehended from the Malays in this district who by intermarriage with the natives have become in some degree naturalized and are indeed much quieter, and more industrious than the Sinhalese or Moors. \citep[61]{Hussainmiya1990}
\end{quote}

Still, it should be noted that during the Dutch period ``the Eastern soldiers, when they embarked for Ceylon, brought their womenfolk along with them. It is not known in what proportion such women from the East-Indies came to the island during the Dutch period. Later the British authorities actively encouraged the foreign Malay recruits to bring along their families.'' \citep[47]{Hussainmiya1990}.


A critical assessment of the extent of and evidence for intermarriage over the whole period of SLM history can be found below in Section \ref{sec:slmbg:RelationtotheMoors}.


\section{British rule}\label{sec:slmbg:HistoryduringBritishrule}
\subsection{Political history}\label{sec:slmbg:British:Political history}
The British were the third European power to rule Ceylon.  They took over from the Dutch in 1796. The Malay soldiers defended fiercely and are even said to have been the only worthwhile enemy of the British \citep[cf.][42,44]{Powell1973}. This earned them the respect of the British. Many of the Malay prisoners of war taken by the British returned to Ceylon shortly after that, serving in the British native army.

\subsubsection{The Kandyan Wars}\label{sec:slmbg:TheKandyanWars}
The first British governor of Ceylon, Governor North was sympathetic to the Malays and valued their military skills highly. He created a regiment of Malays, which changed names several times during its existence: Malay Corps, His Majesty's Malay Regiment, the First Regiment, and finally the Ceylon Rifle Regiment.
At that time, the British regarded the Malays as an exclusive soldier class. Even civilian Malays were registered in order to enlist them later \citep[61]{Hussainmiya1990}.
The Regiment became the central institution of Malay life. Younger boys could serve  as half-pay boys, a name which indicated that they earned exactly half the wage of an older soldier \citep[97]{Hussainmiya1990}, and an invalids' company took care of the ones that had been less fortunate in serving His Majesty. \citet[73]{Hussainmiya1990} indicates that about 75\% of the Malay population served in the Regiment during that period. From 1802 onwards, retired Malays were sent to the South and given land to colonize: first Hambantota, later Kirinda and Palatupana \citep[63]{Hussainmiya1990}.

The British kept on depending on Malays for their supply of native soldiers, because the Sinhalese and Tamils proved to be poor soldiers \citep[24]{Hussainmiya1990}. Most of the able-bodied Malays were recruited into the army, and, as a consequence, there were not many civilian Malays left \citep[73]{Hussainmiya1990}.
But their number could not meet the needs of the British. Recruitment from abroad was encouraged and companies were formed in Hambantota and Tangalle \citep[64f]{Hussainmiya1990}. In 1800, recruitment from Penang in Malaysia was started. This was the first time that Malays from the Peninsula came to Ceylon.  This recruitment policy gave new life to the society of  Malays remaining  in Ceylon \citep[73]{Hussainmiya1990}, with a total of 172 new recruits in 1801 and 1803.

The exiled princes were not recruited but given an allowance because North estimated them important for the social cohesion of the Malay population for the spiritual and cultural leadership they provided \citep[79]{Hussainmiya1990}.

% \begin{quote}
%     Malays were required to carry passes \citet[61]{Hussainmiya1990}
% \end{quote}

When the British attacked Kandy in 1803, many Malays deserted and
joined the Kandyan king fraternizing with the Malays who had
already deserted during Dutch times. \citet[166]{Codrington1926} gives the number of 450 deserters out of a total 700 Malay soldiers. Following this war, the
Regiment only numbered 600 men \citep[72]{Hussainmiya1990}
compared to 1400 at the end of the Dutch reign. Governor North
left the island shortly after in 1805 and was replaced by Governor
Maitland. The latter was aware of the Malay desertion during the
Kandyan war and blamed the British defeat on them. He wanted to
get rid of the Malay soldiers as soon as possible. His plans were to abolish the whole regiment, but he failed even in dissolving  the invalid Malays' company, which he said to be expensive and unserviceable \citep[67,78]{Hussainmiya1990}. Not only the soldiers were affected by his resentment; he also sent back some political exiles, thus depriving the Malays of their cultural leaders \citep[79]{Hussainmiya1990}.


%\begin{figure}
%    \centering
%        \includegraphics{pics/soldierpop}
%    \caption{Development of the numbers of the regiment}
%    \label{fig:soldierpop}
%\end{figure}


%\begin{quote}
%    [Two more companies were added to each of the  three native Bataillons in Ceylon in June 1808] \el ``It is doubtful that there was an actual increase in the number of Malay soldiers in response to this order, for Maitland \el pointed out that the average strength of the Malay soldiers would be only about 600 men for some time to come''\citep[79]{Hussainmiya1990}
%\end{quote}

\begin{sidewaysfigure}
\centering
	\includegraphics[width=\textwidth]{pics/Englishimmigration}
  \caption[Provenance of the Malay immigrants during British rule]{The provenance of the Malay immigrants during   British rule.  The first soldiers came from Jakarta; later Penang and the Straits settlement on the continent became more important. The size of the arrows corresponds to the number of immigrants.}
 \label{fig:provenanceBritish}
\end{sidewaysfigure}

The next lasting governor, Robert Brownrigg,\footnote{Glossing over the brief intermezzo by John Wilson 1811-1812.} held the Malays in high esteem. As a first action, he recruited the 150 Malays that Maitland had not cared to draft. Seeing that the numbers of Malays were running low as a consequence of his predecessor's policy, he started a recruiting mission in Java and recruited 412 soldiers in 1811, accompanied by 214 women and 208 children (see Table \ref{tab:RecruitsFromMalaysia} on page \pageref{tab:RecruitsFromMalaysia}). On the next military campaign against Kandy, he counted on his Malay Regiment to win the war, a hope which was fulfilled \citep[68,81f]{Hussainmiya1990}.

Under Governor Brownrigg, the Regiment further extended the domains in Malay life for which it was responsible. A Royal Military school was founded where the Malays were taught Literary Malay and religion \citep[96]{Hussainmiya1990}.

In 1816, the British protectorate over Indonesia stopped, and
recruiting from Java was no longer possible. The recruiting would
subsequently be done on the Malaysian peninsula.

% \begin{quote}
% In working out the new ratio thus established between Malay soldiers and civilians, it seems likely that after North became Governor, more that 75\% of the local working Malay population had accepted military service as a regular occupation. \citep[73]{Hussainmiya1990}
% \end{quote}


% Kandy prince thama thuukar story h90:71
\subsubsection{Malays in Kandy}\label{sec:slmbg:MalaysinKandy}
The central kingdom of Kandy had been a refuge for Malays
deserting or escaping from the colonial powers ever since the
Portuguese times. Malays fighting for the Sinhalese king Rahju
against the Portuguese in 1587 are mentioned by
\citet[26]{Powell1973}. During Dutch and British rule, more Malays
had escaped the colonial forces and fled to Kandy.  The first
Malay  reported to have done so during Dutch times was a royal
exile who slipped into Kandy in 1742 \citep[69]{Hussainmiya1990}.
In 1765/66 the Dutch General Baron van Eck left behind a large
number of Malay soldiers when retreating to Colombo
\citep[37f]{Powell1973}. Generally, the deserted soldiers and
runaway slaves were held in high esteem by the Kandyan king
\citep[78]{Hussainmiya1990}. The British could not prevent
desertions to Kandy either \citep[67]{Hussainmiya1990}.

The Kandy Malays remained true to their military inclination  in Kandy as well and served as soldiers, guards, and even as bodyguards for the king. After the British-Kandyan wars, the Kandyan king Wickremasinghe was eager to assure their loyalty. But Wickremasinghe lost his temper at the end of his reign \citep[71]{Hussainmiya1990} and the loyalty of the 200--300 Kandy Malays was split in the 1815 war with the British \citep[82f]{Hussainmiya1990}. The Kandyans lost the war, sealing the fate  of the  last Sinhalese kingdom in Lanka \citep[220ff]{Desilva1981}.

% \begin{quote}
% The fate of the Kandyan Malays after the war is not certain. \el Their former leader Assana `Kapitan' became disenchanted with the British after the war and began fomenting rebellion by joining anti-British elements in Kandy. \citep[83]{Hussainmiya1990}
% \end{quote}
% 
% Furthermore, Moors poured into Kandy after the conquest \citep[44]{Ludowyk1966}, possibly leading to a stronger Muslim community in Kandy.

% \begin{quote}
%     There were Malays already fighting for the Kandyan King who convinced their brethren to desert. They were a great strength for the Sinhalese. They were deserters from the Dutch army and runaway slaves \citep[68f]{Hussainmiya1990}
% \end{quote}

% \begin{quote}
%     [anxiety if the Kandyan monarch to placate his Malay subjects''\citep[83]{Hussainmiya1990}
% \end{quote}



% \begin{quote}
% There seems to have been much disgust and dissatisfaction among the 200 or 300 Malay subjects serving the Kandyan Monarch \citep[82]{Hussainmiya1990}
% \end{quote}

\subsubsection{The Ceylon Rifle Regiment}\label{sec:slmbg:TheCeylonRifleRegiment}
After the complete conquest of Ceylon by the British in
1815, the Malay Regiment enjoyed a high status. In 1827, the
Ceylon Rifle Regiment (CRR) was founded. Despite the change of the
name, the conditions of service remained pretty much the same for
the Malay soldiers. But soon after the formation of the CRR, there
was a shortage in recruits from the local Malays
\citep[83f]{Hussainmiya1990}. The Malays seem to have changed
their attitude towards military career. An important aspect was
the reduction of pay. Other reasons were the enlistment for life,
which was not popular, and the service with Sepoys and Kaffirs,
two groups which the Malays held in low esteem. Furthermore, there
were now many possibilities of employment in other professions
due to the growing economy, such as government, police and
plantations. So the military  no longer had a monopoly on the Malay
work force. Failing to attract local Malays, the British again
tried to recruit from abroad and opened a recruiting station in
Penang in 1833  (see Table \ref{tab:RecruitsFromMalaysia} on page \pageref{tab:RecruitsFromMalaysia}). This station proved successful in the beginning, but this would change. The recruiting station shifted to
Singapore in 1840, was stopped in 1842, revived in 1845 and
finally abandoned in 1848 due to its meager success.

The Malays recruited during this period are said to be of ``low quality'' compared to the native Malays and the recruits from Java in 1813 and 1816, and rarely could secure higher ranks \citep[95]{Hussainmiya1990}.
During the time of the CRR, Ceylon saw a small but steady influx from the Malay peninsula of about 50 recruits a year. This influx is thus an order of magnitude lower than the 412 soldiers recruited in 1811 alone.

But Ceylon was pacified and there was not much to do for the CRR.
To make use of the men under arms,  a Malay detachment was sent to
serve in Hong Kong from 1847 to 1854. Service in Hong Kong was not
popular and casualties were high \citep[99]{Hussainmiya1990}. As a
consequence of this, the already low popularity of the CRR with
the Malays fell further: ``By the 1850s it became relatively
difficult to find Malays willing to serve in the CRR''
\citep[88]{Hussainmiya1990}.

In 1848, the Regiment was engaged against a rebellion. This was its last military action. The regiment continued to shrink. In 1849 the Regiment was short of 483 men, in 1851, 22 companies were reduced to 18 companies, and in 1854, to 14, of which 11 were Malay. In 1865, the detachments in Hambantota, Badulla, Kurunegala and Jaffna were withdrawn \citep[101]{Hussainmiya1990}. The police department took over some of the duties from the CRR and many Malays switched over. From 1870 onwards, the Malay element was dominant in the police \citep[102f.]{Hussainmiya1990}.

With no warfare at all in 25 years, it did not seem acceptable to finance the CRR any further and the Regiment was disbanded in 1873, leaving  the remaining 700  Malays to take up new occupations \citep[99]{Hussainmiya1990}. 300 of them are known to have joined the police.
Few discharged soldiers went back to the SEAA. Even after the disbandment of the Regiment in 1873, most Malays stayed in Ceylon \citep[96]{Hussainmiya1990}.
 
% 
% \begin{quote}
%     ``During North's time, the half-boys establishment was fixed at 80, but later the number seems to have been increased from 200 to 300''\citep[97]{Hussainmiya1990}
% \end{quote}


% \begin{quote}
%     [the Malay boys always proved to be the most valuable class of men in the regiment and monopolized the posts of non-commissioned officers]\citep[97]{Hussainmiya1990}
% \end{quote}
% 
% \begin{quote}
% The Malay Regiment had not only become an integral part of the British native army, but had also firmly established itself as a large family unit of Malays bound by strong kinship ties.
% \end{quote}






\subsection{Demographic history}\label{sec:slmbg:DemographichistoryduringBritishRule}

\begin{table}
    \centering
        \begin{tabular}{rrl}
        year & number & commentary\\
        \hline
        1801 & 81     & \parbox[t]{8cm}{
                                         recruits with their families sent by the Raja of Cochin (Malays
                                         employed by the Dutch who had defected to the Raja)}\\
        1803 & 82       &\parbox[t]{8cm}{recruits from Penang, some accompanied by their families
                          and}\\
           & 9      &\parbox[t]{8cm}{boys}\\
      1811 & 412        &\parbox[t]{8cm}{soldiers from Java accompanied by}\\
             & 214    &\parbox[t]{8cm}{women}\\
           & 208    &\parbox[t]{8cm}{children}\\
      1813 & ?      &\parbox[t]{8cm}{recruits from Madura (no number given)}\\
      1816 & 228    &\parbox[t]{8cm}{Javanese from Semarang and Gresek accompanied by their
                        families}\\
      1819 & 102    &\parbox[t]{8cm}{recruits from Penang}\\
      1830 &  50    &\parbox[t]{8cm}{recruits from the Straits of Malacca and Singapore}\\
      1834 & 100    &\parbox[t]{8cm}{recruited in the Straits Settlement}\\
      1835 &  64    &\parbox[t]{8cm}{recruited in the Straits Settlement}\\
      1836 &  50    &\parbox[t]{8cm}{recruited in the Straits Settlement}\\
      1837 &  37    &\parbox[t]{8cm}{recruited in the Straits Settlement}\\
      1838 &  29    &\parbox[t]{8cm}{recruited in the Straits Settlement}\\
      1839 &  37    &\parbox[t]{8cm}{recruited in the Straits Settlement}\\
      1840 &  32    &\parbox[t]{8cm}{recruited in the Straits Settlement}\\
      1841 &  33    &\parbox[t]{8cm}{recruited in the Straits Settlement}\\
      1845-46 &73   &\parbox[t]{8cm}{recruited in the Straits of Malacca and Singapore}\\
      1856-57 &7    &\parbox[t]{8cm}{recruited in Borneo, Pahang, Trengganu, and Kelantan}\\
        \end{tabular}
    \caption[Arrival of recruits during the British Period]{Arrival of recruits during the British Period based on \citep[60ff]{Hussainmiya1987} and \citet[18]{Bichsel}. Note that the total number of Malays recruited from abroad during
North's government is not known.}
    \label{tab:RecruitsFromMalaysia}
\end{table}



\begin{figure}
	\centering
		\includegraphics[width=1.00\textwidth]{pics/Malaypopovertime}
	\caption[Development of the Malay population in Ceylon over time]{The development of the Malay population in Ceylon over time. Data are scarce and unreliable. Data for soldiers are good from 1764 onwards. The other groups are mostly based on two or three years where information is available and linear growth or decrease to/from these years. Furthermore, there is virtually no information about women, so the diagram is to be taken with a huge grain of salt. Sure points are the beginning and end of the exiles population, the short interlude of Kandy Malays and the growth of the free Malay population after 1796 to the detriment of the soldier population.}
	\label{fig:Malaypopovertime}
\end{figure}

By the end of the Dutch rule, the Malays appear ``to have been on a firm footing in the island much so that the British administrators, with very few exceptions, identified and referred to them only as Malays [and not by more precise denominations SN]'' \citep[55]{Hussainmiya1987}.
After a period of consolidation following the British takeover,
the British took the Malays in service and, after all, not much had
changed for them: they continued to serve in the army. The total
number of Malays under arms was 1000 in January of 1800 before the
recruiting from abroad started, then 1200 in 1801 \citep[62]{Hussainmiya1990}. The new recruits were allowed to bring their wives and children (\citet[66]{Hussainmiya1990}, \citet[12]{Hussainmiya1987}). The influx of recruits from abroad
can be seen in Table \ref{tab:RecruitsFromMalaysia}.
The new recruits were the more important as Governor Maitland had
sent back a good deal of the royal exiles in 1808, thus leaving
the Malays in Ceylon without their ``cream'' and ending the days of
aristocratic influence. The newcomers of 1811 made up about a quarter of the Malay population.\footnote{The number
 of Malay soldiers was 1200 in 1801, ca. 1280 if we count the 1803 recruits. Given the near-total recruitment of the Malays through the  Dutch, and the British policy to also accept civilians in  the ranks of the Regiment, we can say that the total number of Malay males combining soldiers and free Malays cannot be much higher than 1600 in 1805 when Gvr. North left. (This does not include women and children.)  Gvr. Maitland had this  number reduce further
through sending the exiles back to Java and not replacing the
casualties. If we assume that 200 exiles were sent back with their
families, then a reasonable guess for the male Malay population in
1811, when the first recruits from Java came, is maybe 1300. Comparing
this to 412 recruits arriving in 1811 with their families, the newcomers made up about a quarter of the Ceylonese Malay population.}

The subsequent recruitments were done in the Malay peninsula and were a lot less massive. The recruits came in a much more trickling way (see Table \ref{tab:RecruitsFromMalaysia}), with  peaks in 1819 and 1834 attaining 100.
%
%There are some records of recruits from abroad having brought their wives and children. This was encouraged by the British \citet[12,52]{Hussainmiya1987}. It is thus to be estimated that the number of immigrants is much higher than the mere number of recruited soldiers. 
After the Regiment had  plunged to 1000 men in 1850, the rules of the regiment were changed and  henceforth  allowed retirement after 15 years of service. Only a few Malays made use of this and went back to Singapore or Penang.
In Colombo and Kandy the Malays were numerous enough to populate entire quarters of the town: Katukalle in Kandy and Slave Island in Colombo \citep[111f]{Hussainmiya1990}.


% \begin{quote}
%     ``Further investigation is necessary to confirm whether the Dutch also followed a similar policy. In any case, the number of women arriving from the East Indies was limited, as a good proportion of 'Malay' soldiers had to find their wives among the local women from the Sinhalese, Tamil or Moor communities. It appears that the Malay Muslims preferred to marry the local Moor women because of their common religious background. A number of such cases of intermarriage between the Malays and the Moors is reported in the Tombos'' \citep[52]{Hussainmiya1987}
% \end{quote}

% \begin{quote}
% [Frederick North] sent special recruiting missions to the East \el{} to bring not only Malay men, but also Malay women and children.\citep[12]{Hussainmiya1987}
% \end{quote}


% \begin{quote}
% ``that men so entering into this Corps will be permitted to take their wives and families with them free from all expense'' (1802) \citep[66]{Hussainmiya1990}
% \end{quote}
 
The disbandment of the regiment did not have a great impact on the composition of the Malay population in Sri Lanka.  Only 80 Malays went back to Singapore \citep[104]{Hussainmiya1990}. But there were no more new recruits from the East bringing fresh blood and new cultural developments to Ceylon.  This lacking cohesion was detrimental to the Malay society. And while the number could be maintained, the degree of internal organization and cohesion was a lot smaller than during the CRR's days. The families had lost their great point of common interest. The new jobs in the police or on the plantations involved work on a smaller scale, in which one would meet a lot less Malays on an everyday basis than before.
Even if many Malays served in the police, it was not an exclusive Malay domain like the Regiment, and Sinhala or Tamil officers would have forced the use of English, Sinhalese or Tamil in the office. Furthermore, the critical mass for maintaining a language might not have been attained by the smaller police stations.

%\begin{quote}
% %   [After disbandment, ~500 Malays in police]\citep[103]{Hussainmiya1990}
%\end{quote}
 
To sum up, during the time of the Regiment, the Malays saw an important
portion of their fellows leave: the aristocrats sent back in 1808.
Governor Maitland did not recruit more Malays, which had the
number of enlisted Malays plunge to 800. Governor Brownrigg wanted a
strong Malay Regiment and perpetrated a massive recruitment, first
from Java, and after the British had lost Java, from Penang and
Singapore. Until 1819, 1016 new recruits joined the ranks of the
Sri Lanka Malays, about evenly from Java and the peninsula.
% The
% Javanese were regarded higher though \citep[93f]{Hussainmiya1990}. These are likely to have
% perturbed the established linguistic setting because of their
% punctuated arrival and their higher prestige.
It was more and more difficult to find recruits in the East. The
number of soldiers decreased. After 25 years without bellicose
activity, the CRR was disbanded. While this did not entail a
back-migration to the SEAA, it stopped the arrival of ``fresh
blood'' to Ceylon and deprived the Malays of an important domain
for using the Malay language.

\subsection{Sociological History}\label{sec:slmbg:SociologicalHistoryundertheBritish}
 
The Sri Lanka Malays can be said to have developed a distinct
identity in the beginning of the 19$^{th}$ century,  different from both
Lankans and Malays in the SEAA \citep[11]{Hussainmiya1987}.
 Their small community was organized around the military service, with the
exiles being cultural and spiritual leaders, who sometimes also
served in the regiment. The exiles formed an
important part of the Malay society, and they ``lent a certain
amount of cohesion, reflected glory and guidance to its
members'' \citep[79]{Hussainmiya1990}, but a good part of them was
sent back in 1808. The military officers then assumed the role of
community leaders. They were well respected in the Regiment, even
among British officers. Until about
1850 the Regiment formed a central part of the Malay society with
every\-thing revolving around it. Security, solidarity and
understanding among the families characterized the life in the
cantonment. The houses were built close to the military
administration center, as was the mosque, where worships were held
in the Malay language. Children could attend the school of the
Regiment and become half-pay boys later before entering the ranks
of the CRR, whereas old people were taken care of by the Invalids'
Company. Marriage within the Regiment was frequent and the
social status of a community member could be evaluated by the
number and prestige of officers attending his wedding \citep{Hussainmiya1987,Hussainmiya1990}. 


 


%\begin{quote}
%   ``A few members of the exiled Royal families seemed to have stayed behind, especially those who married local women. \el the often-claimed princely ancestry of many later day Sri Lankan Malays can rarely be \el substantiated with any credibility''\citep[61]{Hussainmiya1987}
%\end{quote}
 
``The Regiment constituted the centre of Malay-life in Sri Lanka throughout  the British rule during the 19$^{th}$ century'' \citep[93]{Hussainmiya1987}.
The Malays were highly regarded in the regiment, also by the British officers. A career from a half-pay boy to a captain was possible, contrary to European officers, who were taken from the military academy and never from the low ranks \citep[122]{Hussainmiya1990}.
The Regiment was bound by strong kinship ties. Many Malay families depended on the Regiment as their only source of income \citep[78]{Hussainmiya1990} and soldiers joined because of family tradition.
%  The Malay network was tightly knit and the Malays had a certain feeling of superiority regarding the Sepoys and Kaffirs.
The Regiment took care not only of the active soldiers, but also of retired soldiers, who were either given land in the south or cash. Taking land in the south does not seem to have been very popular with the Malays because it would have removed them from the community. Unless the land was close to the cantonment (for instance \em Kampung Pensen \em in Kandy), the Malay preferred the money.
While the Regiment provided cultural and religious support for the Malays, the actual living conditions of the soldiers seem  to have been pretty poor. In the South, things seem to have been worse and soldiers are said to have lived in ``wretched huts'' \cite[115f]{Hussainmiya1990}.

% \begin{quote}
%     [regulations and wages \citep[60]{Hussainmiya1987}]
% \end{quote}


In 1812, the first military school for Malay boys was founded. Their number rose to 5 by 1860 \citep[99]{Hussainmiya1990}. The goal of the school was to provide good education to the ranks of the Malays.  Officers could order their soldiers to attend school. Even if parents were not in the military, they could send their children to this regimental school, which many did. The education provided was excellent and included Malay language, Tamil and English, but no Christian religious education  \citep[96]{Hussainmiya1990}.
\citet[25]{Saldin2001} states that English began to be spoken in Malay homes at that point in time due to economic necessity, although it is unclear on what sources he bases that statement.
In the 19$^{th}$ century, the literacy rate among Malays was the highest among all ethnic group in Ceylon (\citet[48]{Marga1988}, cited in \citet[17]{Bichsel}). There are some numbers of the students attending the regimental schools:

\begin{quote}
In 1835 there were 149 Malay boys studying at military school belonging to the CRR \el in 1873, when the CRR was disbanded it was said that there were 269 half pay boys serving the regiment. \citep[97]{Hussainmiya1990}
\end{quote}

In 1872, shortly before the Regiment was disbanded, the enrolment was 391 \citep[17]{Bichsel}. If we consider the falling number of soldiers in the late 19$^{th}$ century, it is impressive that there were still that many students in the Regimental schools. This must be due to civilian Malays sending their children to the Regimental schools as well.

% \begin{quote}
%     ``Their mode of instruction was intended with a view to make them ``read and write their own and the English languages, so that they might be perfectly capable of translating backwards and forwards the one language into another'' Among the general subjects taught in the regimental school were the Malay, Tamil and English languages.''\citep[96]{Hussainmiya1990}
% \end{quote}
%


The attractiveness of the Regiment began to decline in the 1820s.
After 1850, with low pay and service in Hong Kong making the
Regiment unpopular, the number of free Malays (\em Orang
Priman\em) went up.  
Civilian Malays would engage in the following occupations:
      civilian government departments,
      private business establishments,
      private European plantations,
      civil police department,
      district government agent's establishment,
      coffee and sugar plantation overseers, conductors, security guards, tea-makers,
      servants,
      gardeners,
      rattan weavers,
      clerks,
      peons, and 
      book-keepers  \citep[90]{Hussainmiya1990}. 

Police service proved particularly popular, maybe because of the
similarity to military work. 90\% of the police sergeants in
Colombo were Malay in 1833. After the coffee blight in 1872,
tea cultivation began in Ceylon.   The Malays found work in the
new estates and in the railway construction companies that
connected the highland to Colombo. But contrary to the work in the
Colombo police force, the Malays were more of a minority in the
countryside. For instance, they were outnumbered by the Indian
Coolies in the tea estates \citep[90f]{Hussainmiya1990}.



The British fostered the religion of the Malays and provided a
good infrastructure. There were special Malay mosques next to the
cantonments. This aided the cohesion of the community. Services in
the Moor mosques were  unpopular because the service there was
held in Tamil and Arabic, which did not suit the Malays. This also
sheds doubt on the proposed strong link with the Tamil Moor
community.

\begin{quote}
The soldiers insisted on having mosques situated in the vicinity of their living quarters. Otherwise they had to travel far out of their camps to attend collective worships in the mosques belonging to the Moors. The Malays were at a disadvantage in attending the Moor mosques because the sermons and other religious proceedings were held, not in their mother tongue, Malay, but in Tamil and Arabic. \citep[124]{Hussainmiya1990}
\end{quote}

The mosques did not only serve as a place for religious
meetings. Public meetings could also be held in mosques, also in
Malay language. Mosques were founded  by Malays in
Colombo,
Hambantota,
Kandy,
Kurunegala,
Kinniya,
Gantulawa (Trincomalee),
Galle,
Kalpitiya,
Badulla,
Kirinda
(\citet[124]{Hussainmiya1990}, \citet[449]{Hussein2007}, see Figure \ref{fig:Malaymosques}).

\begin{figure}
    \centering
        \includegraphics[height=.3\textheight]{pics/Malaymosque}
 	\caption{Location of Malay mosques in the 19th century.}
    \label{fig:Malaymosques}
\end{figure}

The cultural leadership of the Regiment was challenged towards the
end of the CRR (1873). In an argument about the position of the
Imam in the Malay mosques, the soldiers could not impose their
candidate, and a certain kind of schism occurred with military
Malays and civilians attending different services
\citep[125]{Hussainmiya1990}.
 
The Malays had regular contact with the SEAA during the 19$^{th}$ century. While they had a distinct identity, they had not forgotten their roots, and the new recruits arriving mainly in the first half of the century assured that the ties would not weaken. This would continue until the death of the last retired soldiers of the Regiment \citep[13]{Hussainmiya1987}. The recruiting missions were often accompanied by Malay officers and    ``some members of such recruiting missions stayed long enough in the peninsula, at times more than two to three years at a stretch, to bring back not only new developments in Malay culture, but also Malay literate texts and manuscripts to be distributed in their community'' \citep[13]{Hussainmiya1987}. Hussainmiya also affirms that the Malays enjoyed visiting Penang, Malacca and Singapore. It is not entirely clear whether he sees this as some kind of tourism, or whether he refers to the recruiting officers.

% \begin{quote}
%     ``even after the disbandment of the Ceylon Rifle Regiment in 1873, which had effectively put a stop to direct links with the peninsula, determined efforts by some members of traditional Malay elites in Sri Lanka, such as Baba Ounus Saldin (b. 1838 - d. 1906) prolonged some form of cultural links with Malayan archipelago until the very late part of the 19$^{th}$ century'' \citep[13]{Hussainmiya1987}
% \end{quote}



% \begin{quote}
% During the period dominated by the traditional Malay elites, that is when a vigourous Malay classical literary tradition was still alive \el{} the Malays in Sri Lanka can be said to have been very much conscious of their places of immediate origin. Thus for examples, Baba Ounus Saldin often took pride in emphasizing that his family had originated from Sumenap in the island of Madura in the East, Similarly, obituary notices inserted in the local press always mentioned the country of origin of the dead persons. \el{} Thus, although the local Malays had been long naturalized in Sri Lanka, the country of their adoption, the memory of their birth  places exercised a strong influence in their national identity with the motherland
%  \citep[13]{Hussainmiya1987}
% \end{quote}

\section{After the Regiment}\label{sec:slmbg:PostregimentalHistory}
 

% \begin{quote}
%     [On the 15$^{th}$ if August 1873 the Regiment was disbanded. There were only 389 soldiers left]\citep[103]{Hussainmiya1990}
% \end{quote}

% \begin{quote}
%     ``the CRR was disbanded in 1873 discharging over 700 Malay soldiers''\citep[99]{Hussainmiya1990}
% \end{quote}




\subsection{Political history}\label{sec:slmbg:postreg:Political history}
The disbandment of the Regiment had a major impact on the Malays.
Many of them now served in the police, but contrary to the
Regiment, the police was organized in small stations all over the
country. They lost a lot of their social cohesion as compared to the
regimental days. Occupations reported by  \citet[10]{Saldin2003} include:
 jail guards; prison overseers;
  firemen;\footnote{\citet[432]{Hussein2007} states that 75\% of the police force, 90\% of the staff of prison services, and 100\% of the Colombo Fire Brigade Quarters were Malays.}
 overseers in the inland tea plantations, which were just being created as a result of
the coffee blight;
and railway workers.
All these jobs had in common that the tight cohesion of the Regiment
was lost and social life was no longer centered around the
barracks. \citet{SmithRH} comments that this disbandment cannot be seen as catastrophic, but it appears that the loss of social cohesion could still be seen as the major change within the social history of the Malays. Hussainmiya illustrates this:

\begin{quote}
    Nevertheless, as for the Malays there was a major difference between serving in the Police Department and serving in the Regiment. As policemen, the Malays were required to serve in the most remote corners in the island, mostly in small units, and therefore their community came to be disbursed in far away places unlike during the regimental days, when they could live in comparatively large numbers in the few military cantonments, which in fact was a prime factor that helped in strengthening the social cohesiveness of the community. \citep[194]{Hussainmiya1990}
\end{quote}


Hussainmiya's object of study being the Regiment, his report ends with the aftermath of the disbandment. We have very little information about the events in the remainder of the 19$^{th}$ century. On a more general scale, the turn of the century saw the rise of the movement for independence as well as the rise of communalism of the different ethnic groups in Sri Lanka \citep{Desilva1981,Hellmann1990,NissanEtAl1990}. Sinhalese nationalism resented Christianity, Westernization and British rule. At the same time, the (Hindu and Christian) Tamils called in their rights, and  the Moors asserted themselves as well \citep[12]{Nuhman2007}. The Malay community, too, formed a political association in Colombo for the first time in 1922 \citep[19]{Saldin2001}, the \em  All Ceylon Malay Association \em (\textsc{acma}), which spread to other parts with Malay population \citep[21]{Hussainmiya1987}.

\begin{quote}
The political-identity crisis for the Malays also brought in its wake a strong desire to revive the community's cultural past which had been in the process of slowly being discarded. \el{}  The Malays' desire to assert their special heritage has been pursued relentlessly since then by various means, by organizing public lectures on their history, by activating special committees to conduct archival research on their past to search and revive their original Malay Jawi, and Javanese script, to publish books in Malay and to conduct special classes in Malay language etc.
\citep[16--18]{Hussainmiya1987}
\end{quote}

Still, more and more mosque services switched from Malay to Tamil \citep[23]{Bichsel}.

To counter being subsumed into the ``Muslim'' group, the Malays demanded representation in  parliament  \citep[15f]{Hussainmiya1987}. They even denied having any links with the Moors at all  \citep[19]{Hussainmiya1987}.
From 1924 to 1965, the Malays had a deputy  in parliament, who was not elected by universal suffrage, but rather appointed by the government to represent the minority of the Malays, strengthening their ethnic cohesion \citep[20]{Hussainmiya1987}. The Malays' representative was TB Jayah, whose legacy is still valued today. TB Jayah was also part of the delegation negotiating the terms of independence with the British, obtained in  1948. After the 1956 election, which replaced the European-oriented \textsc{unp} with the rurally based Sinhala \textsc{mep}, the Sinhala parties used their majority in parliament to vote a number of laws favouring Sinhala. Among these was  the infamous `Sinhala only' law, making Sinhala the sole official language of the country in 1956, which was resented by the other communities and fueled the ethnic tensions \citep[35]{NissanEtAl1990}. As for the Malays, the mandatory use of Sinhala in the education system meant that they had to devote more time to that language than before, when they had mainly functioned in SLM at home and English in education. This caused a sort of chain reaction, where Sinhala pushed English out of the educational systems, and the Malays, who wanted to preserve their superior command of English as compared to the other groups, chose to educate their children in English, which in turn pushed SLM out of that domain \citep{Bichsel,Saldin2001,AnsaldoEtAl2006fel,LimEtAl2007}.

While many Malays had taken up work in the tea estates after the decline of the regiment  \citep[39]{Saldin2003}, in the latter half of the 20$^{th}$ century, many Malays quit the tea estates and moved to Kandy or other Upcountry towns. This might be related to the nationalization of the tea estates in the 1970s, which led to the good jobs being given to Sinhalese \citep[37]{NissanEtAl1990}. From the 1960s onwards, the \textsc{acma} declined in importance, possibly a signal of an identity crisis among the Malays \citep[21]{Hussainmiya1987}.

The ethnic tensions between Tamils and Sinhalese led to the outbreak of open civil war in the 1980s, which in turn led to a rise of ethnic consciousness in the other groups and culminated in the creation of the \textsc{slamac} (Sri Lankan Malay Confederation) on 18 August 1985 \citep[22]{Hussainmiya1987}.\footnote{\citet[22]{Saldin2001} has 1984.} \textsc{slamac} hosted the 2$^{nd}$  Malay world symposium in the same year \citep[40]{Bichsel}. As of today, \textsc{slamac} is defunct \citep[23]{Saldin2001}, and its functions have been taken over by \textsc{coslam} (Conference of Sri Lankan Malays). \textsc{slamac} was founded four years later than the Sri Lanka Muslim Congress (\textsc{slmc}), which is mainly supported by Moors. The Moors were able to turn the \textsc{slmc} into a political party with some influence. This influence led to   recent statistics subsuming both Malays and Moors  under the ethnic [sic] group ``Muslims''. This makes future representation of the Malays in parliament unlikely, since their ethnic group of ``Muslims'' is already represented in parliament by the SLMC, even if all of the members are Moors, whose interests very often diverge from the Malays'. This arrogation of voice of all Muslims by the Moors has led to irritation within the Malay community, who do not feel that the Moors can speak on their behalf, albeit that they are Muslims \citep[24f]{Nuhman2007}. The up-to-now last Malay MP, elected in 1989 (MH Amith), is then indeed not a member of the \textsc{slmc}, but of the   United National Party (\textsc{unp}).

\citet[38]{Bichsel} resumes the political dilemma of the Malays in Sri Lanka as follows:

\begin{quote}
The Malays are in the uncomfortable position of being a minuscule minority  within a much larger and more influential minority [the Moors, S.N.] in a political climate  unfavorable to minority interests. \citep[38]{Bichsel}
\end{quote}

% \begin{quote}
%     ``For the first time the local 'Malay elites', though a miniscule [sic] minority, began to demand their own political rights in Sri Lanka, the country of their adoption, when they felt the opportunity to claim a seat in the legislative council''\citet[15f]{Hussainmiya1987}
% \end{quote}

% \begin{quote}
%     ``The greatest concern of the minority of Malays was that their f and identity was getting over-shadowed by the numerically superior Muslim-Moor group.''\citet[16]{Hussainmiya1987}
% \end{quote}

% \begin{quote}
%     ``The period between 1920 -- 1930 constitutes a watershed in the evolution of ethnic-consciousness among the Malays of Sri Lanka.  They decided to form a political association for the first time \el The motive was to urge the British authorities to concede a Malay seat in the legislative council.''\citet[16]{Hussainmiya1987}
% \end{quote}

% \begin{quote}
%     ``The political-identity crisis for the Malays also brought in its wake a strong desire to revive the community's cultural past which had been in the process of slowly being discarded.''\citet[16f]{Hussainmiya1987}
% \end{quote}

% \begin{quote}
%     ``during the golden Jubilee year of the Colombo Malay Cricket Club in 1922, the Malays decided to widen the scope of a communal association by founding an all Ceylon Malay association with the view to cater to the cultural, economic and religious needs of the community''\citet[16f]{Hussainmiya1987}
% \end{quote}


% \begin{quote}
% The member of the Malay associations vociferously put forward their case for a separate seat for the Malay community It is significant  to note that some of them even went to the extent of disclaiming at this juncture any affinity with the Muslim-Moor community even on grounds of a common religion and the [sic] forgetting their long years of fraternal association. \citep[19]{Hussainmiya1987}
% \end{quote}

% \begin{quote}
%     ``In any case the right to send one of their own men to the legislative council was conceded to the Malays, and from 1924 to 1952 at least one Malay was chosen to represent the community's interests, in 1952, they temporarily lost the Malay representative in the Parliament, but through the six special seats allotted for special interests in the Ceylon Parliament in accordance with the provision of the Soulbury constitution of 1948, a Malay member was nominated to the parliament but only until 1965. Since that year, no Malay member was nominated even for a special seat. The Malays lost this privilege of a nominated member of the parliament with the introduction of a new constitution in 1972.\el The significance of Malay political representation in strengthening their communal identity in Sri Lanka cannot be underestimated.''\citet[20]{Hussainmiya1987}
% \end{quote}

% \begin{quote}
% In the new occupations such as police, prisons and estate services which accommodated a relatively large number of the working Malay population, there was little need to continue education in Malay. In order to be successful in the new careers, the Malays had to learn English and/or one of the local languages, either Sinhala or Tamil. \citep[147]{Hussainmiya1990}
% \end{quote}


% \begin{quote}
%     ``Founded in 1922, the All Ceylon Malay Association (ACMA) had been a principal rallyi[n]g force of the community. [Its] headquarters is situated at Padang, in the heart of the Malay enclave of the Slave island area in Colombo, and its branches spread to other parts of Sri Lanka which have significant Malay population.''\citet[21]{Hussainmiya1987}
% \end{quote}


% \begin{quote}
%     ``special committees were set up [by the ACMA] to collect oral and written literature; also Malay dancing and music were revived.''\citet[21]{Hussainmiya1987}
% \end{quote}

% \begin{quote}
%     ``The loss of importance of the ACMA as the main Malay organisation during the last two decades, also symbolised perhaps the crisis in the ethnicity identity of the Malays.''\citet[21]{Hussainmiya1987}
% \end{quote}

% \begin{quote}
%     ``aggravated demands for separation and self-rule on the part of the Tamils, a largest minority community Sri Lanka [sic], have spurred feverish activities on the part of all the communities to define their respective rights. The birth of a new umbrella Malay organization in 18 August 1985, SLAMAC, or Sri Lankan Malay Confederation (its Malay name being KORAMEL, Konfederasi Rakyat Melayu Langkapuri) can be seen as direct response to new political developments in Sri Lanka''\citet[22]{Hussainmiya1987}
% \end{quote}



\subsection{Social history}\label{sec:slmbg:DemographicpostcolonialHistory}
While Hussainmiya provides ample detail about the sociohistoric conditions in which the Malays evolved in the 19$^{th}$ century, information of comparable quality is not available for the 20$^{th}$ century because the Ceylon Rifle Regiment, Hussainmiya's object of study, had been disbanded by then. In what follows, I will sketch the sociohistorical and demographic history of the Upcountry Malays based on extended discussions with several older members of the Malay communities from Kandy, Gampola, Nawalapitiya, and Badulla. I also draw on similar interviews conducted by \citet{Bichsel}.

The end of the Regiment also meant the end of the regimental schools, and thereby the end of education in Malay.\footnote{A private girls' school was run for some years after the end of the Regiment \citep[102]{Hussainmiya1987}.} The Malays switched to English or Tamil education instead (\citet[88]{Hussainmiya1987}, \citet[14]{Saldin2003}).\footnote{Interestingly, the Moors, who are said to be culturally close to the Malays, refused English education to avoid being proselytized \citep[57]{Nuhman2007}.}
 
Assessments of the importance of this event for the evolution of the language vary, as the following two quotes illustrate.

\begin{quote}
The potentially disruptive effects of the disbandment of the regiment produced its aftermath within almost one generation in the community \el{} [U]nlike in the Regiment when the Malys could live in closely-knit cantonments which reinforced their social and geographical cohesiveness, other sectors which employed the Malays, particularly the police service caused dispersion of the community. Above all, in the new areas of employment, there was no need for the Malays to learn their own language. \citep[155]{Hussainmiya1990}
\end{quote}

\noindent
\begin{quotation}
    However wrenching this may have been for the lives of individual Sri Lankan Malays, [the disbandment of the CRR] can hardly be considered the type of revolutionary change required for rapid contact-induced change. \citep[19]{SmithRH}
\end{quotation}

The disbandment was the most catastrophic event of Malay history in Ceylon.   Whether this has had a decisive impact on the language or not can only be decided if  we compare the shape of the language at two points in time: before ($t_0$) and after ($t_1$) the event. Unfortunately, we have no information at all about the shape of the colloquial language before 1900. This entails that a comparison between the language at the points in time $t_0$ and $t_1$ is impossible. This question simply cannot be decided based on the known material.

Still, sociolinguistic theory suggests that  weakening of the social network incurred by the disbandment of the Regiment probably had a major impact on the Malay language. \citet[124]{MilroyEtAl2003} remind us:

\begin{quote}
	Networks chiefly constituted of strong ties support minority languages resisting institutional pressures to language shift; but when these networks weaken, language shift is likely to take place. \citep[124]{MilroyEtAl2003}
\end{quote}

The generation of Malays born around 1880 (thus the grandparents of the oldest speakers) chiefly worked in the army, law enforcement, the railway, or as overseers on the estates, where Malays formed 10\% of this group \citep[26]{Bichsel} (See Table \ref{tab:MalayOccupations1881} for an overview). In those occupations, there was little need for Malay education, and command of English and a local language like Sinhala or Tamil was more important \citep[147]{Hussainmiya1990}. The children and grandchildren went to mixed schools, where Malays, Sinhalese, Moors, Burghers and Tamils\footnote{Estate Tamils (Indian Tamils) or Ceylon Tamils.} were educated together. Generally, Malays either opted for  Tamil-medium or, most often, for English-medium schools \citep[24]{Bichsel}. Sinhala-medium was not chosen until the middle of the century \citep[27f]{Bichsel}. In one school, in the first grades, classes were taught in Tamil with one subject being taught in English, while in the advanced grades, English was the medium of instruction (K051213nar02). Many schools were run by Irish nuns, where education was in English.\footnote{A curious fact is that in one school, Islamic education was taught by Tamils from Jaffna, who were Hindus. Apparently, this did not pose a problem for either of the two communities.}
% Schools could be as far away as seven miles, requiring two hours of walking to reach them in the morning, and the same amount of time to get back home in the afternoon.
Command of other languages was only necessary for working Malays. Since women very often stayed at home and cared for the children, it is possible that they did not have extended contact to the other ethnic groups at all. \citet[22]{Bichsel} even speculates that at the beginning of the 20$^{th}$ century, there might still have been some monolingual Malay women \citep[also see][25f]{Saldin2001}.

\begin{table}
% use packages: array
\begin{tabular}{llll}
\textbf{Professional} && \textbf{Industrial}&\\
government or local government & 769 & mason & 106 \\ 
physician/surgeon or medical practitioner & 5 & draper/cloth dealer & 6 \\
\textbf{Domestic} &  & tailor/millmen & 38 \\
barber & 8 & butcher/meat salesmen & 21 \\ 
\textbf{Commercial} &   & poultry and egg seller & 1 \\
shopkeeper & 19 & fishmonger & 9 \\ 
boutique keeper & 62 & paddy/grain seller & 3 \\ 
general trader & 83 & baker, bread/cake seller & 97 \\ 
petty trader & 22 & vegetable dealer & 3 \\ 
boatmen & 3 & coconut/copra seller & 2 \\ 
\textbf{Agricultural} &  & coffee picker & 6 \\
estate conductor & 78 & tobacco/cigar maker or dealer & 1 \\ 
coffee garden conductor & 8 & betel/arecanut seller & 2 \\ 
horse/cattle trader & 2 & timber feller & 8 \\ 
fisherman & 10 & salt dealer & 1 \\ 
cultivator & 40 & \textbf{Jewellery and Gem Trade} &   \\
coffee garden cultivator & 8 & goldsmith/silversmith & 1 \\ 
  &   & gem digger & 5
\end{tabular}
\caption[Occupations of Malays in 1881]{Occupations of Malays in 1881 following \citet[51]{Marga1988}, cited in \citet[26]{Bichsel}.}
\label{tab:MalayOccupations1881}
\end{table}



Before independence, the Upcountry Malays had close ties with the Burghers on the estates. During this period, many Malays dropped their traditional name prefixes \citep{Saldin2003}. Malay lifestyle resembled more the Burghers', and was less traditionalist than the lifestyle of the Moors.  This is still the case today with the Malays being much more lenient with regard to the literal following of Islam teachings than the Moors, who tend to be more observant. Examples of this are the wearing of the veil, which is more common among Moor women than among Malay women, and the size of the veil, which tends to be less for Malays, if present at all. Regular attendance of religious services seems to be more important among the Moors than among the Malays. An additional factor which might influence this is that mosques operate in Tamil, which is understood by all Moors, but not by the younger generation of Malays.

Sri Lankan Malays are commonly organized in local organizations (\em Sri Lanka Malay Association of Colombo, Uva Malay Association\em, etc.), which cultivate various aspects of Malay culture (song, dance, sports). Songs and dances are normally adapted from either Sinhalese, Tamil or Malaysian/Indonesian records or videos. Sinhalese and Tamil are translated into Sri Lanka Malay, but Malaysian/Indonesian is not. These organizations seem to follow a British model and protocol with treasurers, patrons, minutes, and elections of boards based on proposition and seconding. Most of these associations are affiliated with the umbrella organization \textsc{slamac} (Sri Lanka Malay Confederation)/\textsc{coslam} (Conference of Sri Lankan Malays).
 
After independence from the UK, language policies favouring Sinhala made more and more Malays choose Sinhala-medium schools  if they could not afford expensive private schools, which continued to operate in English.
With the Malays wanting to keep the comparative advantage they had because of their language skills, many opted for English as the language of the home, to complement Sinhala as the language of education. This pushed Malay out of the equation \citep[29f]{Bichsel}. \citet[76]{Saldin2003} comments:

\begin{quote}
    When we were children, we spoke Malay at home but, spoke  English at school, as that was the medium of instruction. Our children were born to a different world where, Sinhala was the official language  and the medium of instruction. They had to study in Sinhala. \el Sheila insisted that we speak English at home for the benefit of the children \el The Malay families that did not have an academic background in English spoke Malay at first. When the children went to school and learnt in Sinhala, they continued to use Sinhala in their homes.\footnote{Also see \citet{Hussainmiya1990,%}[30 footnote 10],
AnsaldoEtAl2006fel,
LimEtAl2007,
Ansaldo2008genesis}.}
\end{quote}


While in former times, women were not necessarily conversant in English, through the rise of female education in the 1930s, more and more women learned English. This language then supplanted Malay also in the household. The following quotations from Saldin's biography attest to this.


\begin{quote}
    We used to speak in Malay in our family household but, after I married, I rarely had a chance to speak in Malay. Malay was an alien tongue at Taprobane. \citep[76]{Saldin2003}
\end{quote}

\begin{quote}
     Malay was hardly spoken in the homes of the members of the Malay Uniques Club. \citep[83]{Saldin2003}
\end{quote}

\begin{quote}
    [In the Malay association] Most of the members of the older generation knew their Malay but due to various reasons were reluctant to speak it.  \citep[94]{Saldin2003}
\end{quote}


The numerical dominance of the Moors leads to the use of Tamil in the mosques, which many Malays do not understand \citep[86]{Saldin2003}. However, the Wekande Mosque in Slave Islands delivers sermons in Malay on one Friday each month \citep[86,113]{Saldin2003}. Furthermore, there are some Sinhala services now, at least in Kandy.


% In Kandy, Kingswood College was particularly popular among Malays (and Burghers) \citep[43]{Saldin2003}.
Around 1957 Malay was introduced as a school subject, but not for long \citep[118]{Saldin2003}. Since 1993, it is again possible to take Malay as a subject in O/L and A/L exams \citep[29]{Saldin2001}. There is no Malay course as such, but the prospective student is required to work on his own with a manual conceived by BDK Saldin, who also grades the exam.

The current sociolinguistic situation for the Malays is such that the most important language to master is English, which is needed for economic progress. The second most important language is Sinhala, which is the official language, and the majority language in the parts of the country where the Malays live. Knowledge of Tamil is less important, and many younger Malays in Kandy do not know Tamil anymore. Still, Malay children in the smaller Upcountry towns acquire  Tamil as well. Malay, finally, is also not regularly transmitted to the younger generations. The older generations of Malays derive an enormous pride from their quadrilingualism  \citep{AnsaldoEtAl2006fel,LimEtAl2007,Ansaldo2008genesis} and it seems that it is their ability to communicate well with every ethnic group which conveys them their ethnolinguistic identity \citep{Bichsel,Ansaldo2009book}. Frequent allusions are made to this fact: ``we know all four languages'', ``we are the only ones to understand everybody'', etc. This is understandable from the fact that the Malays are a tiny minority in Sri Lanka and can simply not afford particularism. The only way for them to get along is to adapt to as many contexts as possible. While the older generations have very good command of grammar, style and register in all four languages, this is not necessarily the case with the younger generation, who often only have full command of Sinhala and a  local variety of English. Knowledge of Tamil is often absent, and full command of Sri Lanka Malay (which would permit to convey any message) is as well. Still, the Malays have the highest literacy rate in Sri Lanka (93.2\%), second only to the Burghers (98.2\%) \citep[32]{Bichsel}.

In the course of the twentieth century, the main fields of occupation of the Malays did not change, and various jobs in the maintenance of public order like the police, jail guards, firemen or soldiers remained popular. Given that after independence the Sinhalese acceded to the ruling positions in these areas, the Malays' loyalty now resides with them, and support for the actions of the Sri Lankan government and the Sri Lankan Army (both dominated by Sinhalese) in the civil war is widespread.\footnote{I do want to emphasize the complexity of the Sri Lankan society and stress that a simplistic equation government=Sinhala is not warranted. The government is not entirely constituted of Sinhalese, and its actions do not have the support of all Sinhalese.} However, jobs in the government sector were predominantly given to   Sinhalese from the 1960s onwards (\citet[29,33f]{Bichsel},  \citet[37]{NissanEtAl1990}, \citet[55]{OsmanEtAl2008}), a fact which is resented by the Malays, who had had a long tradition as public servants, and the English language skills necessary for that post before the rise of Sinhala nationalism.

In recent years, however,  white-collar activities in the commercial sector seem to be increasing at the expense of the traditional occupations \citep[48]{Saldin2001}. Saldin argues that in these areas, Malays have an edge due to their relatively superior command of English. Also, temporary work in the Middle East is common \citep[34]{Bichsel}, actually since the 1970s \citep{Saldin2001,OsmanEtAl2008}, as is emigration, mainly to Australia and Canada, where emigrant Sri Lanka Malay associations can be found.\footnote{SLAMA in Australia \url{http://www.geocities.com/slama\_web/}, SLAMAT in Toronto.} In the 1990s, there was an initiative for massive emigration to Malaysia, the \em Malay Action Front\em, but it failed to make a noticeable impact \citep[45ff]{Saldin2001}.

The Malays are quite well integrated into the Sri Lankan society and participate in all walks of life up to the topmost strata, including supreme court judges \citep[24]{Nuhman2007}, national sportsmen, ministers, film stars, singers, media presenters, important journalists \citep[53]{OsmanEtAl2008}, or high representatives of Sri Lankan Airlines. This contrasts with the assessment by \citet{Hussainmiya1990} (probably gained in the 1970s) that

\begin{quote}
 there are no noteworthy Malay businessmen and therefore no surprise if one cannot find wealthy people among them. In this respect, the Malays find themselves on the lowest rung of the economic ladder as opposed to other Muslim groups in the island, such as the Tamil-speaking Moors and the traditional business communities of `Borahs' and `Memons'.
\end{quote}

While it is true that the Malays are less inclined to commerce than other parts of Lankan society, there are a number of wealthy Malays in Colombo, but their wealth normally stems from dependent work. This is in line with Bichsel-Stettler's analysis of the Malays' economical situation in the beginning of the 20$^{th}$ century:

\begin{quote}
The Malays were not rich entrepreneurs, nor were they strongly represented in prestigious fields such as law or medicine; nevertheless, as a group they led a solid existence in urban centers. \citep[27]{Bichsel}
\end{quote}

Today, it is possible to marry members of other communities, but the spouses are then required to convert to Islam if they were following another religion before \citep{LimEtAl2006aclcwp}. As a rough estimate, about 10\% of marriages involve a member of the Moor community and 5\% involve a member of the Sinhala Buddhist community. Marriages to either Hindus or Christian are very uncommon. Among the higher strata of society, intermarriage may be more common. \citet[83f]{Saldin2003} remarks that at one stage in the nineties the key office bearers of the Malay Club were children of non-Malays or  were married to non-Malays. This is in opposition to former times, when, according to \citet[64]{Saldin2003}, mixed marriages led to ostracization.

As of today, it appears that the Malay culture is clearly Lankan. Funeral and wedding customs are almost identical to the Moors \citep[1,79]{Saldin2003}, with a combination of Muslim, general South Asian and Western practices,  while the customs relating to a girl's growing up have the greatest similarity with the Sinhalese \citep[35]{Saldin2003}.

%  \citet[17]{Saldin2003} notes that the Malays are superstitious and are fond of making vows to a saint ``for each and every desire''[13].





% \begin{quote}
%     ``In those [former] times the groom rarely saw the bride until he lifted the veil from her face whilst she was seated on the bridal throne, It was his parents and elders who had seen her and deemed her fit enough to be accepted into their family.''\citet[64]{Saldin2003}
% \end{quote}




% 
% \begin{quote}
%     ``\el in Sri Lanka, the sermons in the mosques are delivered in Tamil, which I do not understand \el In the Wekande Mosque in Slave Island however sermons are delivered by Guru Muzni Ameer, in Malay on one Friday in each month and during Ramazan and Haj Festival service.'' \citet[113]{Saldin2003}
% \end{quote}
% 
% \begin{quote}
%     ``\el regardless of the fact that there were Malays whose mother tongue was not Tamil and who had studied in Sinhala.''\citet[86]{Saldin2003}
% \end{quote}



% \begin{quote}
%     Malay mainly in households where English is rarely used, Elites speak more English and less Malay \citet[30 footnote 10]{Hussainmiya1990}
% \end{quote}


\section{Cultural history}\label{sec:slmbg:CulturalHistory}
In the course of history, Malay poets have produced a good number of literary works in Sri Lanka, mostly during the period of the CRR.\footnote{There is no indication of literary activity preceding the Regiment. If there was literary activity, it could not have been widespread and would be confined to religious texts \citep[143]{Hussainmiya1990}.}



%\begin{quote}
%   ``By helping the soldiers to run their own recreation room, which also served as a repository of books, the Regiment can be given the credit for establishing the first Malay library of its kind in the Malay world''\citep[146]{Hussainmiya1990}
%\end{quote}
%club

During the Regiment, artistic work as well as newspapers were
produced. Artistic work comprised adaptions of traditional topoi
as well as the creation of own poems. Genres found in Malay literatures are: religious kitabs; books on
magic, sorcery and divination; catechism and prayer; family trees;
and personal memoirs \citep[135]{Hussainmiya1990}. There is a
penchant towards religious, particularly mystic literature made
for specialists, emphasizing the observance of \em Sharia' \em and the quest for \em ma'rifat \em (gnosis).
The genre of `Syair' enjoyed particular
popularity \citep[139f]{Hussainmiya1990}. Hussainmiya further notes
that  historiography of the SEAA dynasties is completely absent. This can
be understood as an affirmation of an own Sri Lanka Malay identity. On the other
hand, there are two works that are unique to Ceylon: `Hikayat Raden
Bagus Gusti' and `Hikayat Indera Kuraisy', whether they are
originally from Ceylon or just did not survive in the rest of the
Malay speaking world \citep[135f]{Hussainmiya1990}.

Literature was written in High Malay, not in Sri Lanka Malay or other varieties of Malay. The script used was a modified form of Arabic without dots, Jawi script, also called Gundul.\footnote{\em Gundul \em means ``hairless'', and the dots were seen as hairs.} The name \em gundul \em has a Javanese etymology, which might indicate that the exiles from Java had brought literature to Ceylon \citep[134ff]{Hussainmiya1990}.

Sri Lanka Malay literature was not isolated. Grammars and dictionaries were exchanged between Ceylon and the SEAA \citep[145]{Hussainmiya1990}. Despite some localisms, the Sri Lankan Malay literature did not differ much from those works \citep[149]{Hussainmiya1990}.
Two newspapers were issued in Ceylon, \em Alamat Langkapuri \em and  \em Wajah Selong\em. The former was popular with the established Malay population while the latter was preferred by the more recent arrivals \citep[156]{Hussainmiya1990}.

Literary interest was fostered by the British by the creation of a library-like recreation room for Malay officers \citep[146]{Hussainmiya1990}. Malay literature is not only to be read but also to be recited. The Regiment provided a frame for this. With the disbandment of the Regiment, this kind of literary (inter)action became much  more difficult. The poets moved away in search of work and it was difficult to have them meet. By consequence, literary activity stopped in the beginning of the 20$^{th}$ century. The last officers educated in the regimental schools died out, and the younger members of the community did not know how to write Malay. In their new civil occupation, knowledge of English and Sinhala or Tamil was more important. The first literary work in Roman script was published in 1906, but it is of inferior quality  \citep[149]{Hussainmiya1990}, and its appearance was ephemeral. As long as the literature persisted, the Malays were aware of their origins in the SEAA, whereas when the literary tradition died with the disbandment of the Regiment, the links with the SEAA went more and more into oblivion \citep[13f]{Hussainmiya1987}. Towards the turn of the millennium, several Sri Lankan Malays  profited from the time allowed by their retirement to write books about Malay culture and language. These books are written in English. Examples include \citet{Saldin2001,Saldin2003} and \citet{Thaliph2005}.

Literary activity seems to have been confined to Colombo and Kandy \citep[146]{Hussainmiya1990}.  \em Kampung Katukalle \em and \em Kampung Pensen \em in Kandy are mentioned frequently as places of copying the manuscripts.


% After retirement, the officers would still frequent their old
% colleagues in the library room and cultivate literature. This continued beyond the 
% disbandment of the regiment,   but no
% younger Malays joined them, and when they died, the use of Malay as
% a literary language in Ceylon died with them. There were some
% subsequent attempts to produce Malay literature, but they could
% not meet the standards set by the officers.
% \citet[130]{Hussainmiya1990}  sees them as  ``the last bastion of
% the traditional Malay society symbolized in the institutions of
% the CRR''.

As of today, Sri Lanka Malay is mainly a spoken language \citep[cf.][134]{Hussainmiya1990}. Occasionally, some songs or other artistic material are written down in Latin script, Sinhala script, or Tamil script. The regular radio program \em Suara Melayu \em  \citep[established 1957,][28]{Saldin2001}, for instance, contains a popular radio drama, which is enacted by normally three persons reading out the text written in Tamil script. Another example is a text on the life of the prophet Mohammed, found in Kandy, which is written in Sinhala.  On the first Friday of each month, and on some Islamic holidays, the state television broadcasts a Malay program. \citet[47]{Bichsel} reports that letters to the producers of the radio drama are written in Malay using Sinhala, Tamil or Latin script. Starting in  the 1980s, there have been some attempts to set up language classes for the Sri Lankan Malays, which do not seem to have gained wide popularity \citep[39ff]{Saldin2001}. The Malaysian representation in Sri Lanka is very eager to promote Malaysian language and culture among the Malays abroad, while the Indonesian representation shows less enthusiasm in this regard.

There are a number of traditional songs called \em panthongs \em in Sri Lanka.  These used to be accompanied by the violin and  a kind of tambourine called \em rabaana\em. The last person in Kandy to accompany the \em panthongs \em on the violin died some years ago, while the last Kandyan to perform the chant and rabaana died in 2008. Newer Sri Lanka Malay songs are adaptions of Sinhala, Tamil, Indian or English tunes and follow general Indian/Western styles. Some of the Sri Lankan Malays in Kandy have a vast array of modern musical instruments and perform(ed) in bands for weddings and other festivities. Depending on the ethnicity of the bride and groom, they would adapt the set for the evening to Sinhala, Tamil, or English songs, with the occasional Hindi in between.
 
% \begin{quote}
%    ``Malays of Sri Lanka have also had a preoccupation with the mystical aspects of religion. The orthodox texts emphasized the observance of \em Sharia' \em as an external guide to life, while the texts on Sufism urged Malays to engage in time-honoured speculation on man and his place in the universe. The quest for \em ma'rifat \em or gnosis was a strong instinct among the members of the community''
% \end{quote}

% \begin{quote}
%    ``The Malay mystical writings found in Sri Lanka do not appear to exhibit the age-old conflict between pantheistic and orthodox sufistic ideas, which polarised the Malay Sufi thinkers particularly during the late 16$^{th}$ and 17$^{th}$ centuries.''[no pantheistic texts have been discovered]
% \end{quote}


% 
% Cultural authority was exerted by priests and
% officers.\citep[138]{Hussainmiya1990}

% \begin{quote}
%     ``Despite the intrusion of some localists, the \em Syairs \em produced by the writers of the old order \el were in style and language not much different from works written elsewhere in the Malay world''\citep[149]{Hussainmiya1990}
% \end{quote}


% \begin{quote}
%     ``One of its copies made in 1886 by Ahmad Saaly Lye, probably a descendant of our author, has quite a higher quota of `localisms' ''\citep[124]{Hussainmiya1987}
% \end{quote}



\section{Malay settlements}\label{sec:slmbg:Malaysettlements}
Being soldiers, the Malays had always lived close to the centers of colonial power when they were not engaged in military campaigns. The native Lankan population was much more dispersed, and even today, Sri Lanka is not a very much urbanized country.\footnote{Ranked \#181 among 199 surveyed countries, \url{http://www.nationmaster.com/graph-T/peo_per_liv_in_urb_are} retrieved on Sun Jun 28 09:06:52 CEST 2009.} Moreover, the exiles were confined to the centers of colonial power because they had to be under control of the authorities. This urban base distinguishes the Malays from the Sinhalese and Tamil Hindus, and to a lesser degree from the Moors.

During Dutch times, Malay regiments were stationed in the  main coastal towns of Colombo, Galle and Trincomalee  \citep[50]{Hussainmiya1990}.
Kandy was still independent during that time, but had a certain population of deserted Malays.
After the British had vanquished the Dutch, and later the Kandyans, the main Malay settlements were Colombo and Kandy. The first non-military settlement was founded in 1802: Hambantota. Retired Malays were given land to cultivate there. This colony  was followed shortly by Kirinda and Palatupana not far away.

Besides Colombo, Kandy and the South, there were further
detachments in the towns already inhabited by Malays during Dutch
times, and in Badulla, which had been under Kandyan rulership. Following the 1848 rebellion, special detachments
were set up in Kurunegala and Matale. In 1860, there were five
companies in Colombo, three in Kandy, two in Galle, two in Trincomalee,
one in Badulla, and  the Sepoy company in Jaffna. In addition,
detachments of soldiers from Colombo or Kandy or Galle were sent
to Puttalam, Kurunegala, Hambantota and Chilaw. Figure
\ref{fig:Malaypop1860} shows the towns where the Malays were
placed in 1860. In 1865, the Military Commission recommended the
closing of all CRR stations for reasons of economy, except those
of Colombo, Kandy, Galle, and Trincomalee \citep[101]{Hussainmiya1990}. \citet[425ff]{Hussein2007} states that
\begin{quote}
Malay settlements in the other upcountry districts such as Badulla and Bandarawela seem to have originated with the lure of employment in the plantation estates of these areas. These appear to have been principally peopled by disbanded or retired soldiers of the Malay garrison stationed at Kandy. More recently, Malays settlements have arisen from colonization schemes sponsored by the state. The earliest such scheme was Beragama Malay Colony founded in 1930, and this was followed by Bolane Malay Colony off Ambalantota. Recent years have also seen a drift of Malays from Colombo to Wattala, particularly to the relatively new settlements such as Akbar Town in Hunupitiya named after a well known Malay Mas Thajun Akbar. \citep[425f]{Hussein2007}
\end{quote}
The former Malay populations of Kalpitiya and Trincomalee (and possibly Jaffna) seem to have absorbed in the surrounding Moor population \citep[37,426f]{Hussein2007}.

\begin{figure}
    \centering
        \includegraphics[height=0.3\textheight]{pics/Malaycomp}
    \caption{Malay companies or detachments in 1860}
    \label{fig:Malaypop1860}
\end{figure}

According to the censuses of 1981 and 2001, the Malay population concentrates in the population centers given in table \ref{tab:MalayPopCtr}.
\begin{table}
    \centering
        \begin{tabular}{lrrr}
          & Town & Population 1981  & Population 2001\\
        \hline
        1 & Colombo & 20,041 & 19,428\\
        2 & Gampaha & 8,077 & 12,333\\
        3 & Hambantota & 4,380 & 7,385\\
        4 & Kandy & 2,648 & 2,210\\
        5 & Badulla & 1,300 & 1,390\\
        6 & Kurunegala & 1,201 & 1,123\\
        7 & Nuwara Eliya & 1,113 & n/a\\
        \end{tabular}
        \caption[Population centers with Malay population]{Population centers with Malay population. Census data stem from \citet{Bichsel} for 1981 and from \citet[423]{Hussein2007} for 2001. Notable is the gain of population in Gampaha district, where many Malay families from Colombo moved to. The population loss in Colombo was countered by immigration from other areas. The rise of Malay population in Hambantota is surprising and cannot be explained here.}
        \label{tab:MalayPopCtr}
\end{table}

The variety of SLM described in this grammar is spoken in the central Upcountry town of Kandy. Kandy was the last seat of the Sinhalese kings, and the last kingdom to surrender to the British. It is the cultural centre of the Sinhalese and houses the Temple of the Tooth containing the most important Buddhist relic in Sri Lanka, a tooth of the Buddha. Kandy is the second largest city of Sri Lanka, with a population of about 100,000. Historical population figures given by the Municipal Council\footnote{\url{http://www.kandycity.org/kandycity-geogra.html}, retrieved on Wed Mar 25 11:34:49 CET 2009.} are 2500 for the early 19$^{th}$ century and 16,881 for 1871.

\begin{figure}
	\centering
		\includegraphics[width=\textwidth]{pics/Kandymap}
	\caption[Map of Kandy town]{Map of Kandy town. Total around 250 families.}
	\label{fig:Kandymap}
\end{figure}

Kandy lies at an altitude of 500m, at the foot of the central tea cultivating region. The main road (and railway) accessing that region leaves from Peradeniya, at about 5 km distance from Kandy. The first Malays to arrive in Kandy were soldiers deserted from the colonial armies. This happened during Dutch times. It is not entirely clear whether the Portuguese as well had  already Malays under their service (see \ref{sec:slmbg:MalaysinKandy}). The Malay population centered around the quarters of Bogambara and Katukalle, where there were also Malay mosques (which still bear that name today). Especially the Old Peradeniya Road leading through Katukalle and Mulgampola contained many Malay households. However, this area has developed economically and many new shops have opened there. This has led many Malays to sell their property there and to move to more rural areas. The roads leaving from Peradeniya Road southbound into the hill area through Mulgampola, Rosawatta and Heerassagala also contain many Malay settlements. But today, Malays are not confined to one area and can be found interspersed in all quarters of Kandy. The Malay population in Kandy is about 2,500, making up  2.5\% of the total population. Even in quarters with relatively many Malays, the percentage would not be higher than perhaps 5\%, making the Malays a not  very visible minority. In fact, there are other inhabitants of Kandy who are not aware of the existence of a Malay minority in their town at all. Figure \ref{fig:Kandymap} gives an overview of the repartition of the Malay population in Kandy.


The main upcountry towns of Kandy, Nawalpitiya, Nuwara Eliya, Hatton and Badulla were by no means isolated, and there were frequent contacts among inhabitants of different towns, which can be seen by the very common marriages involving partners from different towns. Marrying into a family of another town was a common practice (and still is today). This might have  helped preserve Upcountry Malay as a coherent variety with little internal diversification. As an interesting side effect, the idiolectal diversity we find in the Upcountry seems to be tied not to region, but to family lines. For instance, one speaker from Kandy uses \phonet{a\E P\dz u:PuP} for the present tense from \trs{aràduuduk}{is.sitting} even in careful speech, which is normally pronounced \phonet{ar\E \dz u:\dz uk}. This makes more sense if one knows that she was born in Badulla, but that her father was born in Hambantota in the South, where such contractions seem to be more common. The rest of her idiolect does not differ much from other idiolects from Kandy, but this one feature betrays the origin of her family in the South two generations earlier.


\begin{figure}
 \centering
 \includegraphics[height=0.3\textheight]{pics/Kandy-network}
 % Kandy-network.jpg: 478x600 pixel, 362dpi, 3.35x4.21 cm, bb=0 0 95 119
 \caption[Strength of contacts between Kandy and other towns]{Impressionistic strength of contacts between Kandy and other towns, and between Badulla and the South}
 \label{fig:Kandynetwork}
\end{figure}

Contact with Colombo was also very frequent (at least after independence), and most Malays from Kandy have relatives in Colombo. Other than that, one speaker from Kandy had her mother from Negombo, to the North of Colombo, and  another's mother is from Puttalam. The wide-spread kinship ties of the Malays prove that in the last three generations, the Malay network covered all of the Upcountry and the western coast with Colombo and Negombo. It is less common to have relatives in the South (Hambantota, Kirinda); this community seems to form a separate network \citep{AnsaldoEtAl2006fel,LimEtAl2007}. There are some Malays in Galle, who still speak Malay, but who seem to be quite isolated from the rest. Figure \ref{fig:Kandynetwork} gives an impressionistic illustration of the perceived strength of ties. Further, more sophisticated demographic research must show whether the postulated strengths of ties hold up to statistical analyses.


The Malays in Kandy have long been organized in cultural associations, the oldest one of which is the Kandy Malay association. Recent years have seen a number of schisms and splits among the Kandy Malays, leading to the foundation of the \em Hill Country Malay Club \em and the \em Hill Country Malay Association\em. The Malays in Kandy take part in island-wide cultural activities, such as singing, dancing, and sport events. In distinction to certain strata of Colombo society \citep{AnsaldoEtAl2006fel,LimEtAl2006aclcwp}, there seems to be little interest in teaching Malaysian Malay in Kandy, and attempts to speak Malaysian Malay are occasionally derided as \em Thuan thuan language\em, imitating the reduplicating pattern found in Malaysia, but not in Sri Lanka.

While in the surrounding Upcountry towns, the Malay language seems to be regularly transmitted to the younger generation, this seems to be less the case in Kandy, where English and Sinhala have made inroads into the domain of child care.

% \begin{table}
% 	\centering
% 		\begin{tabular}{lll}
% 		A & Sudhumpola	& 25	\\
% 		B &	Katuakalle 	& 20	\\
% 		C &	Mulgampola  &	20\\
% 		D &	Siyambalapitiya & 10	\\
% 		E & Heerassagala 	&	 30, scattered\\
% 		F &	Bowala, Dangole& 10	\\
% 		G &	Peradeniya &	2\\
% 		H &	Rosawattha& 30	\\
% 		I &	Nuwara Dodawala& 6	\\
% 		J &	Deiyannewella&	10\\
% 		K &	Bahivavakonda &	2\\
% 		L &	Asgiriya &	3\\
% 		M &	Malnaiyama&	10\\
% 		N &	Buwelikada & 10	\\
% 		O &	Wewelpitiya & 0	\\
% 		P &	Ampitiya &	5\\
% 		Q &	Lewella & 10	\\
% 		R &	Mapanawatura &	1\\
% 		S &	Mawilmada & 6	\\
% 		T	&	Kathugasthota & 20 \\
% 		U & Kahalla 5 &	\\
% 		V &	Watupuluwa &	0\\
% 		V &	Malwattha &	0\\
% 		V &	Kotugodella &	0\\
% 		V &	Kandy Town &	15\\
% 		V &	Hanthane &	3\\
% 		\end{tabular}
% 	\caption{List of Malay families in quarters of Kandy}
% 	\label{tab:ListOfMalayFamiliesInQuartersOfKandy}
% \end{table}
% 
% \begin{table}
% 	\centering
% 		\begin{tabular}{ll}
% 		 \textbf{Matale}	& 50	\\
% 		 \textbf{Kurunegala}	& 170	\\
% 		 Nawalapitiya	& 150	\\
% 		 Gampola	& 200	\\
% 		 Kegalle	& 30	\\
% 		 Nuwara Eliya	& 20	\\
% 		 \textbf{Badulla} 	& 175	\\
% 		 \textbf{Bandarawela}	& 50	\\
% 		 \textbf{Hatton}	& 25 \\
% 		 Ratnapura	& 20	\\
% 		 Balangoda	& 20 \\
% 		 Passara	& 15 \\
% 		 Badalkumbara	& 10	\\
% 		 \textbf{Hambantota}	& 200	\\
% 		 \textbf{Kirinda}	& 50	\\
% 		 Galle	& 60\\\citep[86,113]{Saldin2003}
% 		 Puttalam	& 100	\\
% 		 Chilaw	& 20	\\
% 		 Anuradhapura	& 25 \\
% 		 Mineeriya & 25	\\
% 		 Mihinthale	& 10 \\
% 		 Trincomalee	& 75	\\
% 		 Batticaloa & 10	\\
% 		 Matara	& 20	\\
% 		 Polonnaruwa & 20 \\
% 		 \textbf{Colombo} & 1000\\
% 		\end{tabular}
% 	\caption[List of Malay families in Sri Lankan Towns]{List of Malay families in other Sri Lankan Towns. Boldface indicate good preservation.}
% 	\label{tab:ListOfMalayFamiliesInOtherSriLankanTowns}
% \end{table}




% %"`During the Dambadeniya period (1220--1293 A.D.)Javanese Forces headed bx Chandrabhaanu who professed to be a Buddhist, had made an attempt to invade Sri Lanka.  Thi sattemot was successfully repulsed bz the Sinhalas. [...] The beginnings of Malay/Javanese linguistics influence on the Sinhala language could therefore be traced from about the beginning of the Polonnaruva period (1098--1234 A.D.)"'\citep[97]{Sannagala}


\section{Relation to the Moors}\label{sec:slmbg:RelationtotheMoors}

In all academic studies about SLM, with the exception of
\citet{Jayasuriya2002} and Ansaldo's work, a strong emphasis on the contact between
the Malays and the Moors\footnote{Some Moors resent the term `Moor' for their ethnic group and prefer `Conakar/S\=onahar'.
 `Moor' goes back to the Portuguese who designated Arabs as `Moros', while `Conakar' goes back to Tamils who designed first Greeks and then Arabs that way \citep[2]{Hussein2007}, so neither of those is a native term \citep[45]{Nuhman2007}. In the course of history, there has been a heated debated among this community whether one should identify as `Moor' or `Muslim' \citep[42ff]{Nuhman2007}. I follow academic usage in employing the term `Moor', which furthermore has the advantage of being less ambiguous than `Muslim'.}
is found. This has been called the `Tamil bias' \citep{Ansaldo2005ms,Ansaldo2008genesis}. The proponents of a close link between Malays and Moors argue that the two minorities shared Muslim faith and were thus natural allies in a predominantly Buddhist and Hindu country under Christian colonial reign. This is also used to explain the linguistic shape of SLM,
because  Moor women married to Malay men would have acquired an
imperfect Malay that they passed on to their children. This
imperfect acquisition would be at the base of modern SLM.
\citet{Ansaldo2005ms,Ansaldo2008genesis} criticizes that the historical evidence provided for this scenario is not up to scientific standards, and that oral history of the Sri Lankan Malays suggests that extensive intermarriage with the Moors did not happen.\footnote{Ian Smith (p.c.) remarks that Malays have a political interest in being seen as separate from the Moors (see Section \ref{sec:slmbg:postreg:Political history} above), and that, therefore, they might be downplaying their Moor heritage.}

On analyzing the references in the literature one finds that the claim of close contact to the Moors is mainly based on \citet{Hussainmiya1987, Hussainmiya1990}. The theory does not lack plausibility and is supported by some historical sources. But  Hussainmiya often argues on intuitive grounds and the sources for certain of his claims are not given or vague. He himself  (1987:45) says that ``the details of the nature of mutual contacts between these people cannot be documented.'' Let's take a close look at what he presents.

The argument is based on:

\begin{itemize}
    \item the gain in Islamic knowledge
    \item the Moors helped a Malay prisoner
    \item Christopher Schweitzer's report
    \item the Dutch \em tombos \em and Sri Lankan \em Kadutams \em
    \item a cultural comparison with Malays in South Africa
\end{itemize}

\citet[43]{Hussainmiya1990} argues that Malays in Ceylon must have had contact with great Islamic scholars. Thus, a certain exile named Radin Adipati Natakusuma was banned to Ceylon in 1743. When he returned to Jogyakarta in 1758, he was made chief of the religious officials there. A similar case is that of Wirakusuma, who became the leader of a religious group on returning. Other exiles are said to have become
pupils of the Islamic teachers Sayyid Musa Ngidrus and Ibrahim Asmara in Ceylon. This is seen as evidence for contact between the
exiles and Islamic scholars in Ceylon.  Whether these scholars were Moors is a different question. \citet{Wieringa1997} informs us that Ibrahim Asmara was an Islamic apostle from Champa (today's Vietnam), and as such it is unlikely that he was a Sri Lankan Moor.  It remains furthermore
questionable whether the normal soldiers would have undertaken thorough religious studies in Ceylon.

The second piece of evidence are minutes from the Dutch Political Council Colombo, suspecting the Moors of aiding exiled Malays in communicating with each other.

The third piece of evidence is a report of a German employee of the \textsc{VOC}, Christopher Schweitzer, from 1680, taken as evidence for Malay-Moor intermarriage by \citet{SmithEtAl2006cll}.

\begin{quote}
    Their language is Ambonese, but the majority also speak Malay, Sinhala, Portuguese and Dutch. \el{} The wives, who in part are Ambonese, in part Sinhalese and Malabarian may not say anything [against the stripping of their ornaments]. \citep[106]{Schweitzer1931} [my translation]\footnote{\el ihre Sprach ist Ambonesisch können aber meistentheils ach Malley- Singules- Portuges- und Holländisch \el die Weiber, welche theils Ambonesin, Singulesin und Malabarin, dörffen nichts darzusagen.}
\end{quote}
 


It is difficult to find evidence for contact between the Malays
and the Moors in this passage. There are of course the Malabarian
wives, with \em Malabarian \em being an old word for  ``person speaking a South Indian language'' (mostly Tamil) \citep[cf.][92]{Rogers1990}.
But nothing is said about their religion, so Hindu Tamils cannot be excluded \citep[cf.][68]{Hussainmiya1987}. There are also Sinhalese and Ambonese wives, and the soldiers are not even mentioned to speak Tamil,\footnote{Interestingly,
 \citet[45]{Hussainmiya1990} affirms that ``soldiers also spoke Sinhala and Tamil.'' How he finds this in Schweitzer's text is not clear to me.}
only Malay, Sinhala and the colonial tongues. If one has to rely on this fragment of text, then it surely does not support a close contact between Malay soldiers and the Moors, because this would have resulted in knowledge of Tamil rather than Sinhala.

The fourth piece of evidence are the Dutch \textit{tombos},  particularly the head tombo available at the Sri Lanka National Archive 1/3758, in which  a number of cases of intermarriage are said to be reported. Unfortunately, Hussainmiya does not give the numbers, so it is difficult to evaluate this claim. He (1987:25) also cites Lankan marriage registers (\textit{Kadutams}) that are in his possession and that also list intermarriages. These registers could not be consulted for this book.

\citet{Ansaldo2008genesis}  checked these claims against the photocopies of the Sri Lankan National Archive stored in the Dutch Nationaal Archief in the Hague. He reports that the tombos are inconclusive at best:
 

\begin{quote}
1. The records for the period up to 1796 are damaged by water, making significant
parts of the entries illegible. The most revealing information for identification
here are the signatures of the parties. There is however hardly any information of
ethnic group, which makes it difficult to identify Malay/Indonesian and Moors
given that both groups share the practice of adopting Arabic names. In a
particularly interesting section in the tombos dedicated to mixed marriages (cf.
Hussainmiya 1987), only five of 238 entries clearly refer to individuals of
Javanese origin: of these, two records refer to Javanese-Moor marriages, one to a
Javanese-Javanese marriage, and the remaining two are unclear.\\
2. The following period until 1919, albeit under British rule and therefore less
interesting for our claim, shows a more structured archiving system where
indication of race is given. Where legible, this reveals still a majority of Western
marriages, a growing number of marriages between Eurasians and Burghers
(locally born of Dutch/Western heritage), and between Burghers. There are two
clear entries involving Malays, one married to a Eurasian (between 1867-1897),
and one to a Burgher (1885-1897). From 1897 onwards, race is clearly specified;
of 196 entries only one is Malay. \citep{Ansaldo2008genesis}
\end{quote}

Given the little intelligence which could be gathered, \citet{Ansaldo2008genesis} conducted some pilot interviews with Colombo Malays, who professed not being aware of intensive Moor intermarriage in their ancestry. The time-depth, and breadth of these interviews is limited of course, but they can be complemented by native amateur genealogy. \citet{ArfinBurah} has compiled a book about his genealogy, tracing his origins back to the Javanese exiles. This book contains 303 instances of marriages where both partners are mentioned. One of the partners is in the lineage of the exiled princes and is therefore Malay. The group membership of the other spouse varies. While the ethnicities of the spouses are not mentioned, their names often clearly indicate their background. In Sri Lanka, names can indicate which group the bearer belongs to. European first names are indicative of Burghers or other Christians, e.g \em Coralie Nadine Caspersz \em (Burgher) or \em Wilson Gunaratna \em (Sinhalese Christian). Sinhalese last names (like \em Gunaratna\em) indicate that the bearer is either a Sinhalese Buddhist or a Sinhalese Christian. Tamil last names (like \em Arulambalam\em) can either point to Tamil Hindus or Tamil Christians. Many general purpose Muslim names exist which are borne by Moors and Malays alike \citep[cf.][]{Hussainmiya1987, Ansaldo2008genesis}, e.g. \em Fatima\em, but some names are borne exclusively by Malays.\footnote{\citet[436]{Hussein2007} contains a list of common Malay names. Also see \citet[8]{Hussainmiya1987}} Cases in point are the first names \em Tuan \em and \em Gnei \em or the last names \em Dole \em or \em Sourjah\em. I am not aware of any names which are exclusively borne by Moors. Finally, there are some names in the list which cannot be classified in this schema, very often because they are not complete (e.g. \em Mina\em). Table \ref{tab:slmbg:marriage} breaks down the spouses of the descendants of the exiles. A list of the marriages retrieved is found in the appendix (p. \pageref{sec:marriages}).\\

\begin{table}[h]
\begin{tabular}{rcccccc}
	& Malay & general Muslim & Sinhalese 	& Tamil & Burgher 	& Unclear\\
absolute& 189	& 73		& 14		& 3	& 9		& 15\\
relative& 62.4\%& 24.1\%	& 4.6\%		& 1.0\%	& 3.0\%		& 4.6\% \\
relative w/o unclear cases&65.6\% & 25.3\%&4.8\%&  1.0\%   	& 3.1\%		& / \\
\end{tabular}
\caption[Provenience of spouses of Malays]{Provenience of spouses of Malays in \citet{ArfinBurah}. Total number of marriages: 303}
\label{tab:slmbg:marriage}
\end{table}

The column `general Muslim' mentioned above must be partitioned between Malays and Moors. How this split should be done is unclear, but we can set an upper limit to intermarriage with Moors: No more than a quarter of all contracted marriages involved a Moor. This holds true under the assumption that all instances classified as `general Muslim' actually refer to Moors. At least 2/3 of all Malay marriages did not involve a partner of another ethnic group. The upper bound for in-group marriages is 91.3\%, adding up `Malay', `general Muslim', and `Unclear'. Intermarriage with the Buddhist and Christian community is present but not common; both in comparable proportion of about 4\%. Intermarriage with the Tamil Hindu community is quasi non-existent. Marriage with non-Muslims is found in about 9\% of the cases, but the names given to children of such unions indicate that those families followed Muslim (and indeed Malay) naming practices for their offspring.

The genealogy of the exiles treated in \citet{ArfinBurah} spans about 200 years. It  starts with the exiles in the 18$^{th}$ century and goes through various positions in the Regiment, ending in civil professions in the 20$^{th}$ century. It appears that the offspring of the exiles touches all important domains for Malay culture, and it is unlikely that the marriage pattern followed by the individuals surveyed here is completely unrepresentative of the community as a whole. It seems fair to assume that between 10-15\% of all marriages involved a Moor partner.\footnote{This percentage seems to be valid for the present day as well. \citet[39]{Bichsel} reports that among the   Malays she interviewed, there had been 12 cases of intermarriage with Moors, either by them (22 consultants) or by their parents (44, if none of the consultants where siblings). The ratio is thus 12/66=18.2\%. However, the pattern does not seem to be uniform: \citet[417]{Nuhman2007} reports that in 1956 and 1960 more Malays married outside their community than within. \citet{Ansaldo2008genesis}, on the other hand, reports that ``of approximately 50 families interviewed in total, only two revealed genealogies including Moor-Malay
intermarriage.'' More demographic research is thus needed.} While this reflects the privileged positions of the Moors as compared to the other groups, it is not obvious that this rate of intermarriage would have had a significant effect on the evolution of the language.

Based on \citet{ArfinBurah}, the extent of intercommunal marriage totals about 20\%. As far as the regiment Malays are concerned, the percentage might have been even lower. The British Major Skinner ``attributed the financial difficulties of Malays to continuous inter-marriage in the regiment, and the resultant growth of `swarming connection' of soldiers'' \citep[119]{Hussainmiya1990}. The use of `intermarriage' in this context is not about marrying \em outside \em the regiment, but rather \em inside \em the regiment, i.e. with other Malays.


The fifth piece of evidence is of a rather speculative nature and compares the fate of the Malays sent to South Africa with the Sri
Lanka Malays. The Malays in South Africa have completely lost their original language and culture, which is attributed to the
absence of a coreligionist community such as the Moors in Lanka, which would support them. From the fact that the Malay culture has
not been lost in Sri Lanka,  \citet[53]{Hussainmiya1990} concludes that this must be due to the Muslim community of the Moors. While such an analysis is possible, it has to be mentioned that the disappearance of the Malay element in South Africa has surely more causes than just the
presence or absence of another minority, and that even the presence of a strong Moor community in Ceylon has not prevented Malay culture from being lost in the 20$^{th}$ and 21$^{st}$ centuries. Such a monocausal explanation should be handled with care.

To sum up, we have  the  transfer of Islamic knowledge to some exiled princes; Dutch suspicions;  dubious attested intermarriages in the \em tombos\em; Schweitzer's report, which casts doubt on the relation between Malay soldiers and Moors; and a speculation about South African Malays, which is, well, a speculation. This is far from being solid evidence for a close Malay-Moor relationship or linguistic influence from Tamil. Future research should be aware of the shakiness of these assumptions and take a more sceptical approach, like \citet{Ansaldo2005ms}.

After discussing the points that are said to be evidence for a close relation between Malays and Moors, I would like to briefly shed light on some more facts which make pervasive intermarriage unlikely: At least in the Upcountry, many Malays still have   distinctively South East Asian features, and can easily be told apart from Sinhalese, Tamils, and Moors, which look more South Asian \citep[416]{Hussein2007}. If intermarriage had been that rampant, we would expect that the South East Asian features would have been submerged by now. While South East Asian physical features are easy to be found in the Upcountry, this does not mean that all Malays in the Upcountry look like Indonesians. The spectrum stretches from prototypically Indonesian to prototypically Lankan, with all combinations in between. My point is not that no intermarriage took place, only that it was less pervasive than assumed in the literature. This is reflected in conflicting testimonials from foreigners having come to Sri Lanka at different points in time. \citet[115]{Percival1803} remarks:

\begin{quote}
Although the Malays intermarry with the Moors and other castes particularly in Ceylon and by this means acquire a much darker colour than is natural to a Malay; still their characteristic features are so striking predominant, that they cannot be mistaken.
\end{quote}

Over a century later, a visitor from abroad noted

\begin{quote}
Meeting Ceylon Malays there one cannot help noticing that some of them have the features of Javanese while others look like Malays and their personal names incline to both Javanese and Malay.
  (Said 1926 \nocite{Said1926}, cited in Hussainmiya 1987:34)
\end{quote}

Again half a century later, the former Malaysian Prime Minister Tunku Abdul Rahman remarked in the Malaysian newspaper \em Star \em in June 1981:

\begin{quote}
 They look more like Indians \el{} than Malays and their language is strongly influenced by the Indian dialect. \el{} [A]fter generations of intermarriage, it is hard to pick one from the other Malays, or the Moors, except when they themselves announce their racial identity. \citep[7]{Hussainmiya1987}
\end{quote}

These quotes are of course impressionistic, but show that different people had different impressions about Sri Lankan Malay appearances. If anything can be drawn from those quotes, it is that the Malays became more Indian-like over the course of time. While in the early 19$^{th}$ century, their ``characteristic features are so striking predominant'', and in the early 20$^{th}$ century they ``have the features of Javanese'', in the late 20$^{th}$ century ``they look more like Indians''. These quotes do not support the view that heavy intermarriage was going on in the early days. However, it is most likely that these three different people each met only a tiny subset of Sri Lankan Malays, and that their reports are valid for that subset only. Hence, they do not reflect a global appreciation, and cannot be used as exact sources. As stated above, even today the Sri Lankan Malay community comprises everything from very Indian to very Indonesian, and depending on the Malay one meets, his ethnic identity is easy or impossible to ascertain. It is likely that Percival met Malays of a more Indonesian type, while Tunku met Malays of a more Indian type, but this should not lead us to hasty generalizations.


% Another piece of evidence which speaks against too close association with the Moors is the fact that in Kandy, the Moor community is much less present than in Colombo. While in Colombo, about a quarter of the population is Moor, in Kandy this is much less. Interestingly, the dialectal differences between Kandy and Colombo are slim. Now, there are two possible ways to explain this: either there were more Moors in Kandy in former times, or the linguistic change SLM underwent is not caused by the Moors.

At least in Kandy, the relations between the Malays and the Sinhalese are better than with the Moors. Many Malays do not speak Tamil anymore, while this is still the language preferred by most Moors.\footnote{Some Moors speak Sinhala (\citet[62]{Nuhman2007}, \citet[46]{Hussein2007}).} The Malays are traditionally loyal to the powers that be (\citet[52]{OsmanEtAl2008}, \citet[48]{Saldin2001}), first the colonial powers, then the Sinhalese \citep[cf.][35]{Bichsel}. Since the Moors were not associated with these powers, they are not held in very high regard. This former association with the colonial powers has also led the Malays to adopt a more `Western' lifestyle \citep[cf.][41]{Saldin2001}. All senior Malays asked about their youth affirmed that they used to hang out with Burghers, leading a `happy-go-lucky'-lifestyle \citep[27]{Hussainmiya1990}, which is unusual with Moors.\footnote{A penchant for gambling is already noted in early reports, e.g. \citet[106]{Schweitzer1931} or \citet[131]{Pieris1918}. The literary activities of the Malays may add a more intellectual overtone to this lifestyle, but this would not make them resemble the Moors more. To the contrary, Moors were known as a ``backward community''  \citep[95]{Hussainmiya1987}, unlikely to engage in literary pastime. Furthermore, Malays had the highest literacy among all ethnic groups in Sri Lanka, while Moors had the lowest \citep[48]{Marga1988}, cited in \citet[17]{Bichsel}. As for the present day conflation of Malays and Moors in government statistics, \citet[36]{Bichsel} remarks: ``[A] comparison of the statistics given for the Moors with that for the Malays in the tables in this paper shows clearly how greatly the communities differ.''} Today, most Malays have a very liberal interpretation of Islam, while the Moors tend to be more conservative. This difference in lifestyle again makes overwhelming intermarriage unlikely. \citet[35]{Saldin2003} reports that as far as education is concerned, Malays seem to be closer to Sinhala practices.

After screening the evidence, let us now turn to the plausibility of early contact with the Moors.
If we imagine the situation in the early days of the Malays in Ceylon, would contact between Malays and Moors be plausible?  Would exiles be interested in meeting Islamic scholars? This cannot be excluded, but it does not seem likely that theology was one of their major preoccupations. Islamic scholarship was even less of an argument for the soldiers, convicts and slaves, whose interest in advanced theological discussion can be estimated to have been rather low. The latter groups had, on the other hand, interest in finding a Muslim wife, if they had not brought their family \citep[47]{Hussainmiya1990}. However, intermarriage to non-Muslims is already present in the early days, as shown in \citet{ArfinBurah}.

A further problem for contact between Malays and Moors could be that the colonial powers of the Portuguese and the Dutch tried to contain Muslim influence. In 1626, the Portuguese expelled the Moors \citep[84ff]{DeSilva1972},\footnote{Also see \citet[72f]{Vlekke1943} for the motivations of the Portuguese.} who would settle around Batticaloa, then under Kandyan rule \citep[113]{Codrington1926}.
The Dutch  restricted their share of coastal trade, and forced many  Moors into agricultural life in the East of Ceylon \citep[42]{Hussainmiya1990}.

The Malays, on the other hand, were always close to the centers of colonial administration \citep{Ansaldo2009book},  either because they had to be guarded, in the case of the exiles, or because they were the guards of the colonial power in the case of the soldiers.

After the middle of the 18$^{th}$ century, the percentage of free Malays grew \citep[49]{Hussainmiya1990}, and free Malay could easily enter in contact with the Moors. But the Malays liked to stay close to  their community,\footnote{\citet{Ansaldo2009book}  even argues that \em endo\em gamy was a defining feature of the SLM community at that point in time, rather that \em exo\em gamy as claimed by \citet{SmithEtAl2004} and \citet{SmithEtAl2006cll}.} and the Dutch had just pushed the Moors away from the colonial cities (Batticaloa is an exception to this, and maybe Trincomalee).\footnote{It is interesting that the British later would refrain from stationing a Malay Regiment in Batticaloa.}

Starting in 1783 the Malays had their own mosques, which continued to provide services in Malay afterwards.

\begin{quote}
    As the Malay community grew in numbers, and with the patronage extended by the colonial rulers, their religious needs were fulfilled according to Malay traditions. As the Malays claimed, special mosques for their congregations were set up in the Malay-majority townships, while they were ministered by special Malay \textit{Khatibs} or priests (as styled by the Malays). The Malays also had their own religious \textit{Kitabs} and legal texts written in [t]he Malay language.  \citep[19]{Hussainmiya1987}
\end{quote}


Religion was clearly a domain where the Malays affirmed their identity. Of course, the Moors were fellow Muslims, but it seems that the mosque was not so much a place of interethnic contact but a place of affirmation of group identity \citep{Hussainmiya1987}. Influence through religion could have taken place before 1783 or after 1930, when Tamil began to replace Malay as the language of the mosque \citep[23]{Bichsel}.

%\begin{quote}
%   ``the details of the nature of mutual contacts between [the Malays and the Moors] cannot be documented''\citep[44]{Hussainmiya1987}
%\end{quote}
%
%
%Hussainmiya says that Malays cannot be identified according to their culture, so basically, a Malay is a person that identifies herself as Malay.

%\begin{quote}
%   ``Schwitser did not mention about the religious background of the Amboinese, and it may be that they were not followers of Islam, which explains partly why only Sinhalese and Malabar (Tamil) wives are mentioned in this case. Or it may be that the `Moorish' women were included in the racial terms of `Malabaris'.''\citep[68, footnote 68]{Hussainmiya1987}
%\end{quote}
 



% \begin{quote}
%     ``Intermarriage with other races took place among the Malays [during the time of the Unique Malay Club] \el [at one stage in the nineties the key office bearers of the Uniques were children of non-Malays or, were married to non-Malays]''\citet[83f]{Saldin2003}
% \end{quote}

% \begin{quote}
%     ``In Malaysia and Indonesia it is the groom who gives the dowry, unlike in Sri Lanka, where the bride provides the dowry.''\citet[2]{Saldin2003}
% \end{quote}
% 
% 
% \begin{quotation}
%     ``There are no restrictions on interethnic marriages in Sri Lanka, and there has been intermarriage between the Sinhalese and the Sri Lankan Malays''\citet[9]{Jayasuriya2002}
% \end{quotation}


\section{Language issues}\label{sec:slmbg:LanguageIssues}
After looking at the sociohistorical development of the SLM community, it is now possible to investigate some questions of a more linguistic nature. The first such question is the origin of SLM. Other questions concern the development of the language in Ceylon, neighbouring languages, and the relative importance of different contact languages over time.

\subsection{Sources in the SEAA}\label{sec:slmbg:SourcesintheSEAA}
What are the linguistic ancestors of Sri Lanka Malay? It is clear that the Sri Lankan Malays ultimately came from South East Asia, but which languages or dialects from which islands contributed to its genesis? People from many different islands arrived in Sri Lanka as soldiers after having spent a variable period in Batavia/Jakarta before  sailing off to a battlefield. During the stay in and around Batavia, this soldier population began to acquire a distinct identity, distinct from the Javanese and Sundanese  of Java \citep[174]{Vlekke1943}.\footnote{\citet[99]{Raben2000} actually supports the idea that group identity in Indonesian societies of the time  was generally \em not \em created through ethnicity, but rather through ``personal attachments and charismatic power.''} \citet[14]{SmithRH} argues that  no single mother tongue could impose itself, and the lingua franca of the SEAA, a variety of Malay, was chosen for interethnic communication. % \citep[also cf.][174]{Vlekke1943}.

This variety has been given various names, among which Vehicular Malay \citep{SmithRH}, Trade Malay \citep{Ansaldo2005ms} and Bazaar Malay need to be mentioned.\footnote{It should be noted that Bazaar Malay can also refer to other varieties \citep{Paauw2003}.} Like most trade languages, Vehicular/Trade Malay was not monolithic and would be adapted to the communicative needs of speakers of different ethnic groups in different ports, with different influences from native languages. The sociolinguistic situation of Batavia/Jakarta is traditionally very complex. \citet[132]{Vlekke1943} states that ``in seventeenth century Batavian \el [i]n some respects, the confusion of languages was like that of Babylon''.
 \citet{AdelaarEtAl1996} list the following relevant varieties for the Jakarta region: \em Java Malay, Java Port Malay, Jakarta Malay (Betawi), Jakartan Indonesian, \em and \em  Tangsi Malay\em. It is beyond the scope of this study to investigate the sociohistorical developments of these varieties, but see \citet{Grijns1991} or \citet{Paauw2004,Paauw2008phd} for some discussion.

This complexity suggests that there were many input languages which participated in the formation of SLM. The
different homelands of the Malays arriving have contributed to a
greater or lesser extent to SLM as we know it today. \citet{Bakker1996stuf,Bakker2006} claims that  ``it is fairly well known what Malay looked like  when it was imported to the island by the Dutch''. This opinion is maybe a bit optimistic given that the exact nature of the language (or languages) that the new immigrants spoke is far from clear to the remaining researchers. The total absence of written records from that period does not make the task any easier.
 What can be done is to look at
several features of present day SLM and trace them back to Malay
varieties of the SEAA. This endeavour has been undertaken by
\citet{Adelaar1991}.
He finds that Moluccan, Bazaar Malay, Baba
Malay and Jakartanese have contributed most to SLM, besides the notwithstanding influence of Tamil and Sinhala (Table \ref{tab:slmbg:Adelaar1991}). \citet{Paauw2004} critically reassesses the data and analysis by Adelaar and comes to somewhat different results. He finds that many of the features of SLM are found in several varieties of Vehicular Malay (VM), so that they cannot be attributed to any single variety. Furthermore, he traces more of the features to adstrate influence from Tamil, and discards most of the Moluccan influence, as well as all of Baba and Bazaar influence.


\begin{table}
\begin{tabular}{rp{4cm}p{4cm}p{4cm}p{2cm}}
	& Standard Malay & Sri Lanka Malay & Adelaar &  Paauw\\
\hline
1) &  *h> 		& h*h> Ø & Moluccan, Bazaar, Baba, Jakartanese &  general VM \\
2) &  		& retroflex series & (Tamil and Sinhalese )&  Tamil \\
3) &  		& (contrastive) vowel length & (Tamil and Sinhalese )&  Tamil \\
4) &  		& (contrastive) consonant gemination & (Tamil and Sinhalese )&  Tamil \\
5) &  m/n/\ng/ 	& \ng	&Moluccan&  *Tamil \\
6) &  \E	& i/u or \E{} varying with i/u 	& Moluccan &  *Tamil \\
7) &  retention of most of the inherited morphology
		& loss of most of the inherited morphology & Moluccan, Bazaar, Baba [...] &  general VM \\
8) &  locative preposition +noun phrase
		& noun phrase + linker + locative preposition & Bazaar&   Tamil \\
9) &  noun + determiner
		& determiner + noun &  Moluccan, Bazaar, Baba, Jakartanese&  general VM \\
10) &  possessed + possessor
		& possessor + linker + possessed &  Moluccan, Bazaar, Baba, elsewhere&  general VM \\
11) &  noun + adjective
		& adjective + noun &  (Tamil and Sinhalese)&  Tamil \\
12) &  prepositions
		& postpositions & (Tamil and Sinhalese)&  Tamil \\
13) &  subject-verb-object
		& subject-object-verb & (Tamil and Sinhalese)&  Tamil \\
14) &  ada denoting existence of a noun
		& ada/\E r\E: progressive aspect of a verb & Moluccan, Bazaar, Baba&  *Tamil \\
15) &  		& negators: t\E r/tra &Moluccan, Baba& Moluccan\\ 
16) &  full tense-mood-aspect adverbials &  full and reduced tense-mood-aspect adverbials & Moluccan, Baba&  general VM \\
17) &  plural personal pronouns are independent lexemes
		& plural personal pronouns are historically compound forms with *ora\ng{} & Moluccan, Bazaar, Baba&  general VM \\
18) &  		& 1$^{st}$ and 2$^{nd}$ personal pronouns borrowed from Hokkien Chinese & Bazaar, Baba, Jakartanese&  Jakarta Malay \\
19) &  		& plural marker pada 	& Jakartanese& Jakarta Malay \\
\end{tabular}
\caption[Origin of SLM features according to \citet{Adelaar1991} and \citet{Paauw2004}]{Features distinguishing Sri Lanka Malay from Standard Malay, and their origins according to \citet{Adelaar1991} and \citet{Paauw2004}. Features which Paauw argues to be misanalyzed by Adelaar are marked with an asterisk *.}
\label{tab:slmbg:Adelaar1991}
\end{table}
 
\citet{Paauw2004} further undertook an etymological analysis of the lexicon.
His findings are presented in table
\ref{tab:originsOfTheSLMLexicon}. In this table, origin \textit{Sri Lanka} means all words that joined the vocabulary after the language had reached Ceylon, and includes Tamil, Sinhalese, Portuguese, Dutch and English words. \textit{Common Malay} indicates words understood in every part of the Malay world, whereas \textit{Indonesian} words mostly come from Betawi or Javanese and are not understood outside Indonesia. \textit{Malaysia}, finally, indicates words from the peninsula.

\begin{table}
    \centering
        \begin{tabular}{l|l}
                        origin      & percentage\\
                        \hline
                        Sri Lanka &       10,0\% \\
                        Common Malay &    55,7\% \\
                        Indonesia &       11,7\% \\
                        Malaysia &        0,6\% \\
                        Loan Words &      21,7 \% \\
        \end{tabular}
    \caption{Origins of the SLM lexicon according to \citet{Paauw2004}}
    \label{tab:originsOfTheSLMLexicon}
\end{table}

\citet{Paauw2004} concludes that ``the Malay variety most likely responsible for the lexification of Sri Lanka Malay was a variety originating in Jakarta, and that Moluccan Malay and the Malay of the Malay peninsula also contributed, in a more limited manner.''
Two features which Adelaar and Paauw have not addressed actually corroborate this analysis: SLM has words with schwa in the final syllable (e.g. \phontrs{\dentt a:n@m}{plant}), typical of Jakarta \citep{Adelaar1985}, and SLM drops schwa in antepenultimate position (\phontrs{cri:\dentt a}{story} compare Std. Malay \em c\E rita\em). This is also common in Jakarta (Uri Tadmor, p.c.). Phonological influence from Malaysia is noted by \citet{Paauw2004} in the lowering of high vowels in the last syllable as compared to Indonesian, e.g. \phontrs{ma:sok}{enter}, compare Indonesian \em masuk\em, although \citet[99]{Paauw2008phd} states that in Eastern contact varieties of Malay this lowering can also be found.

% While lexical influence from Malaysia is surely limited, the vowel lowering is very regular and pervasive. The possibility that the language had already stabilized when the first immigrants from Malaysia arrived (1819) gains some support from the lack of lexical influence, but is weakened by the phonological impact these immigrants had. Further research will be needed to ascertain the relative importance of Malaysian vocabulary against phonology.
 

\subsection{Three theories of genesis}\label{sec:slmbg:Threetheoriesofgenesis}
It is a fact that Sri Lanka Malay as spoken today is strikingly different from any other Malay variety. It must have undergone radical language change in the last 350 years \citep{AnsaldoEtAl2009age}. There are competing theories as to how this language change came about, and by what it was triggered. The idea defended by Hussainmiya, Smith and Paauw in their publications \citep{SmithRH, SmithEtAl2004, SmithEtAl2006cll} is that SLM is the product of intermarriage between Malay men and Moor wives, who spoke Tamil. The Tamil speakers would then try to acquire the language of their soldier husbands,   very much like the slave populations in the Caribbean tried to acquire a European language. The language learners would transfer grammatical structures of their original language into the target language. This means that a strong Tamil substrate would remain in the  variety spoken by the wives, which would then be passed on to the children, and by and by generalize to the speech of the whole community. This scenario of genesis is then very similar to the genesis scenario proposed for  Sri Lanka Creole Portuguese \citep{Smith1979}, where the indigenous Tamil population of Batticaloa tried to acquire Portuguese and imposed Tamil grammatical structures on that language. This theory of genesis fits well with the Relexification Theory \citep{Muysken1981,Lefebvre2001} of Creole genesis, where we find a combination of the grammar of the substrate with the lexicon of the superstrate.  In Lefebvre's account of Haitian Creole, the West-African language Gbe would be the substrate and French the superstrate; for SLM, Tamil would be the substrate and the language of the immigrants the superstrate.

A somewhat different position is taken by \citet{Bakker2000convergence,Bakker2000rapid,Bakker2006}, who argues that the variety spoken in Sri Lanka was not heavily influenced by the local languages until the end of the 19$^{th}$ century, citing  documents written by military officers which comply with the general Malaysian/Indonesian standard. Only after the disbandment of the regiment in 1873 would there be heavy language change toward   Tamil (and possibly Sinhala), which would take place within one generation.

\citet{Ansaldo2005ms,Ansaldo2008genesis,Ansaldo2009book} finally proposes that the ecology of Sri Lanka Malay has more in common with settings of metatypy as described by \citet{Ross1996,Ross1997,Ross2007}. SLM would not be the result of intermarriage and  relexification, as proposed by Smith et al. Rather, SLM would be the product of the immigrants' having to handle three languages at the same time: Malay, Sinhala, and Tamil. Given that the grammatical markup of Sinhala and Tamil is very similar, and to lessen the cognitive burden of handling too many different constructions, SLM speakers would gradually emulate the Lankan structures they use everyday in Sinhala and Tamil in their own language as well. SLM would be the in-group language (\em esoteric \em in Ross's terms), which is influenced by the language(s) of wider communication (\em exoteric\em), in this case Sinhala and Tamil. This is analogous to the grammatical change of Takia (\em esoteric\em) towards Waskia (\em exoteric\em) in the Madang province in Papua New Guinea, as described by \citet{Ross1996}.

All these three theories of genesis are tied in one way or another to the historical facts, but they make different assumptions about the historical sociolinguistic setting. Some of these assumptions are explicit and some are implicit, some are supported by evidence, while others are not. In the following sections, I will discuss the socio-historical setting which is assumed for each of the three theories. This will be followed by an explanation of the mechanism which would have led to the change of the language.  I will then discuss the internal plausibility of the setting and mechanism assumed, before I turn to the  historical evidence  presented by the supporters of the particular theory.  This evidence is sociological on the one hand, and linguistic on the other. In the expository parts, I will not engage in criticism of the proposed theories. In the sections after the exposition, both linguistic and sociological evidence presented will be  discussed, critically evaluated, and weighed against other pieces of evidence not necessarily discussed by the supporters. Finally, every theory will be checked for its place in current theorizing on language contact. Does it comply with generally contact linguistic theory, or does it go against some basic assumptions and principles in the field?

For expository reasons, I will start with Bakker,  then turn to the Tamil substrate hypothesis as advanced by Hussainmiya, Smith and Paauw, and discuss metatypy as proposed by Ansaldo in the end.

\subsubsection{The rapid conversion hypothesis}
This argument was the first to be published \citep{Bakker1995nl}. It involves two components: SLM has converged towards Tamil (and possibly Sinhala), and this change happened rapidly. At first, this hypothesis argued that change took place between 1870 and 1906, but later versions \citep{Bakker2006} are more lenient as to the precise start and end dates of the rapid change. In the earlier versions of this hypothesis, convergence was the mechanism at work, while in later work, metatypy is also mentioned as a possibility. In this section, I will limit myself to the discussion of convergence. Metatypy will be discussed more extensively in Section \ref{sec:slmbg:metatypy}.

\paragraph{Argument}
Bakker first mentions Sri Lanka Malay in \citet[17f]{Bakker1995nl}, but the discussion is too brief to really distill any claims about its genesis (besides the analogy to Sri Lankan Portuguese). In \citet{Bakker2000convergence,Bakker2000rapid}, he develops that

\begin{quote}
[Sri Lanka Portuguese and Sri Lanka Malay] were documented in the early part of the last [=19$^{th}$ SN] century and in both cases it is clear that the languages were creolized at that point in time: they lost all person inflection, almost all derivation, and all morphological irregularities. Instead, these languages appeared to have preverbal markers for tense, mood and aspect, a system which may be considered typical for creole languages (Holm 1988)\nocite{Holm1988}. Word order is rather rigid SVO and the languages are typologically close to isolating. In short, they look like a prototypical creole. This, however, is only true for older stages of these languages.
\end{quote}

He then goes on to argue that modern SLP and modern SLM have become more complex and have departed from that creole typology. Bakker argues that ``the whole process took place in one or two generations.''\footnote{also in \citet[153]{Bakker2006}.} 
In \citet[607]{Bakker2000rapid}, he concludes

\begin{quote}
In the case of SLM we must \el{} assume a few decades at most for the radical changes, as earlier documents show a more ordinary form of Malay.

Thus, `convergence intertwining' may change a language structurally and typologically in just a few decades from isolating to agglutinative.
\end{quote}

\citet[152]{Bakker2006} gives some more indications as to the sources of his claims about earlier stages of the language: ``early sources (from 1806 or 1820 to the 1930s) show a form of Malay with little grammatical influence from Tamil.'' Furthermore, he claims that \citet{Hussainmiya1987} contains some instances of Tamilized features in early texts, which he discounts as minor.

Bakker's conclusion is somewhat puzzling:

\begin{quote}
The two main opinions on the matter of \em when \em the convergence happened both have to accept that it is a very rapid process. If it happened soon after arrival, \el{} the convergence or metatypy must have happened within one or a few generations. Alternatively, if it happened in the first half of the twentieth century \el{}, it must also have happened within one or two generations.
\end{quote}

The third possibility, that the change happened between arrival and the second half of the 20$^{th}$ century, as a continuous process, is not addressed.

\paragraph{Sociolinguistic setting assumed}
Bakker's hypothesis  makes few sociolinguistic assumptions. The main and most important assumption is that Tamil was present as an important adstrate during the whole period, while Sinhala influence is also possible \citep[46]{Bakker1995nl}.

\paragraph{Mechanism}
Bakker assumes convergence, or metatypy as defined in \citet{Foley1986} and \citet{Ross1996, Ross2001}.

\paragraph{Plausibility} If the premisses are as Bakker suggests, the developmental scenario he sketches is plausible, although not everyone will agree that language change can happen that quickly.  \citet{Ansaldo2005ms,Ansaldo2008genesis} and \citet{AnsaldoEtAl2007dcctr}, for instance, argue that rapid change is not a possibility for the restructuring of the case system we find in SLM. They argue that a restructuring of a whole functional system (like case) cannot be accommodated under the rapid conversion hypothesis and must necessarily be gradual, requiring more time than allowed for by Bakker.

\paragraph{Historical evidence}
Bakker refers to Hussainmiya's work for the historical evidence he needs, but does not indicate what parts, chapters or pages contain that evidence. It is likely that he refers to the \em Syair Kisahnya Khabar Orang Wolenter Benggali \em \citep[86]{Hussainmiya1987}, from 1820.\footnote{The nature of the earlier documents cited by Bakker, I have not been able to ascertain. Bakker (p.c.) affirms that he based his work exclusively on Hussainmiya. \citet[86ff]{Hussainmiya1987} contains a list of all early texts. None of them predates 1820.} The Syair should show that the Malays used a creole-like language in the beginning of the 19$^{th}$ century. As pointed out by \citet[178]{SmithEtAl2006cll}, this text is in Literary Malay, which, while sharing some features with Creoles (no person inflection, SVO word order), is generally not considered a creole. The linguistic composition of this text can be explained by the British having brought Malay teachers from the Peninsula, who taught Literary Malay to the officers, and is surely not a result of creolization in Ceylon.   \citet{SmithEtAl2006cll} argue that during the period of the Regiment, diglossia existed, with Literary Malay as the high variety and Sri Lanka Malay as the low variety, mirroring the general sociolinguistic pattern found in South Asia in general, and in Sinhala and Tamil in particular.\footnote{Whether
 the restricted use of Literary Malays by some officers can be called `diglossia' is open to debate, but it is certain that Literary Malay was not a language of widespread use in 19$^{th}$ century Ceylon.}
This is further corroborated by the fact that newspapers edited in Ceylon had subscribers in the SEAA \citep[89]{Hussainmiya1987}, who apparently could understand the language of the publication \citep[24]{Saldin2001}.

Bakker argues that the language emerged rapidly. He gives two possibilities for the rapid emergence: either shortly after arrival, or in the 20$^{th}$ century. For the `early' hypothesis, he does not present any historical evidence. He argues that if  diglossia involving `traditional' and `Tamilized' varieties existed in the 19$^{th}$ century, then the Tamilized variety must have developed very quickly, within one or two generations. It is not entirely clear to me why this would be necessary. For diglossia to be present in 1850, any time between 1650 and 1840 would be possible as a `start date' for the Tamilized variety, and more than two generations could be encompassed. As for the `late' hypothesis, he argues that until the middle of the 19$^{th}$ century, the Malays' language was very similar to other Malay varieties. He bases this claim on the existence of texts from that period which do not show Lankan features. Other texts he refers to are from the beginning of the 20$^{th}$ century, and show Lankan features. Thus, the language change must  have happened between these two points in time. Note that for the `early' hypothesis, Bakker assumes diglossia at 1850, while for the `late' hypothesis, he does not. Depending on the hypothesis to defend, his interpretation of the sociohistorical setting in the 1850s is different.

Bakker seems to have an `all-or-nothing' approach to diglossia in the 19$^{th}$ century. One variety must be High Malay, the other one must be either SLM as we know it today, or `Proto-SLM'. However, the possibility of gradual development of SLM cannot be ruled out. It is possible that `Early Low SLM', spoken during Dutch times, gradually evolved into `Modern Low SLM', spoken today, through the intermediate stage of `Middle Low SLM', spoken during British times. This variety would then be in a diglossic situation with High Malay, used in the British Regiment schools and the literary works. In order to argue for rapid conversion, one has to show that either `Early Low SLM'=`Middle Low SLM' with rapid conversion to `Late Low SLM'. Or one has to show that `Middle Low SLM'=`Late Low SLM', with rapid conversion taking place before, ideally within one or two generations. No proof for either of these equations has been brought forward.

\paragraph{Linguistic evidence}
Bakker shows that SLM shares many structures with Tamil, and also mentions similarities with Sinhala in passing. These structures back up his claim of convergence. Bakker does not use linguistic evidence to prove his claim of rapid development.

\paragraph{Counter-evidence}
Since there is no real evidence for Bakker's theory, counter-evidence to refute his claims is not necessary, strictly speaking. Counter-evidence against Bakker's claims can be divided into three domains: 1) against convergence, 2) against early rapid change, and  3) against late rapid change. As for 1), the claim that SLM grammatical structure is close to Tamil structure\footnote{And Sinhala structure for that matter.} seems pretty much uncontroversial, although some non-Tamil features remain in SLM \citep{Slomanson2006cll}. As for 2), since there are no documents available, it is very difficult to make any claims about the structure of SLM preceding 1800. 3), finally, is the only domain where true counter-evidence can be found. \citet{SmithEtAl2006cll} point out that the documents from the middle of the 19$^{th}$ century Bakker refers to are written in Literary Malay by officers who had had schooling in that variety. They argue that at that point in time, there was a diglossic situation, with SLM as the low variety and Literary Malay as the high variety. All literary production from that period is in Literary Malay and cannot be used to make any claims about the structure of the colloquial language. Another argument against the `late rapid change' hypothesis is the absence of strong dialectal differences in SLM.
If the claim is that SLM changed after the disbandment of the regiment which entailed catastrophic changes, we would expect that these changes would only affect the speakers who witnessed them. Speakers in remote places like Hambantota (settled in 1802), with little contact to the other communities, should not change their language \citep[cf.][]{SmithRH}. We would thus expect a linguistic difference between the `catastrophic communities' and the others. While there is some dialectal difference between Hambantota and the other dialects, it is mostly phonological. In any case, it does not point towards the Colombo community having had a significantly different development from the Kirinda community. If we find the same linguistic structures in Colombo and Hambantota, it makes no sense to attribute the structure of the Colombo dialect to some social changes which could never have affected  Hambantota.
 
\paragraph{Place in contact language theory}
The mechanisms of convergence and metatypy are well-established in contact linguistics, although the speed of development Bakker assumes is not a common feature of theories of language contact, and indeed metatypy is generally argued to be a gradual process \citep[127]{Ross2007}.

\subsubsection{The Tamil substrate hypothesis}
This hypothesis is implicitly present in \citet{Hussainmiya1990} and worked out in detail by \citet{SmithEtAl2004, Paauw2004} and  \citet{SmithEtAl2006cll}. It assumes that the development of SLM was analogous to the development of Sri Lanka Portuguese, and analogous to the development of Creole languages in other parts of the world. Speakers of a substrate language (Tamil in this case) try to acquire a superstrate language (Malay), but transpose the grammatical structures of their native language onto the language to be acquired.
 

\paragraph{Argument}
The Tamil substrate hypothesis states that ``SLM developed as a result of intermarriage between Malay men and Moor women,   vehicular Malay  thus served as its lexifier and Sri Lanka Muslim Tamil (SLMT) as its principal  substrate language'' \citep{SmithEtAl2004}.
This is explicated in more detail in the following two passages.

\begin{quote}
From their arrival, [the Malays] associated closely with the established Tamil-speaking Sri Lanka `Moor' community, with whom they  shared the Muslim religion \el. As the soldiers were single males for the most part, a high degree of intermarriage with the Moor community was inevitable. Through this interaction between the two communities, a new language arose with Tamilized structure and with a Malay lexicon. \citep[160]{SmithEtAl2006cll}
\end{quote}


\begin{quote}
From the earliest days in Sri Lanka, `Malays' has a close relationship with the Tamil-speaking `Moors' \el. Intermarriage was first noted by a German observer, Christopher Schweitzer \el, who also commented on the multilingualism of the Malay soldiers \citep[49--50]{Hussainmiya1987}, and was confirmed by Dutch \em thombos \em (administrative records) of the eighteenth century  \citep[52]{Hussainmiya1987}. An indigenized variety of Malay must have arisen among the children of these marriages, whose fathers spoke VM but may have even been familiar with Tamil through their interactions with the Moor community, and whose mothers spoke Tamil natively and must have attempted to learn Malays without the benefit of instruction. Tamil-speaking mothers must also have tried to make Malay the language of the home; otherwise it would be difficult to imagine the language being maintained by subsequent generations.
 \citep[176f]{SmithEtAl2006cll}
\end{quote}


\paragraph{Sociolinguistic setting assumed}
For this hypothesis, three sociological assumptions are fundamental: a) there was heavy intermarriage between Malays and Moors; b) Tamil was the substrate and Malay the superstrate, i.e. Malay had more prestige than Tamil; and  c) the children of bilingual parents acquired the non-native variety of the mother, rather than the native variety of the father. Furthermore, Smith and colleagues discard Sinhala as an important factor. Why this is done is not entirely clear. Presumably it is  because the recombination of the two components grammar and lexicon, which is the basis of the relexification model, is difficult to implement if the number of involved languages is not two.

\paragraph{Mechanism}
Although they do not state so explicitly, Smith and colleagues seem to adhere to a substratist school of creole studies. The idea is that grammatical structure from a substrate is carried over into the superstrate when the latter is learned. The new language is thus a recombination of substrate grammar and superstrate lexicon. Formalizations of this approach can be found for instance in \citet{Muysken1981} or \citet{Lefebvre2001}.

\paragraph{Plausibility}
The substratist/relexification theory is well-tested in many contact settings, and the argument as such is plausible if the assumptions are correct. However, the statement ``Tamil-speaking mothers must also have tried to make Malay the language of the home; otherwise it would be difficult to imagine the language being maintained by subsequent generations'' \citep[177]{SmithEtAl2006cll} is not completely plausible, as \citet{Ansaldo2008genesis} notes. He says that ``there [would] have been [no]  plausible reason for
Tamils/Sinhalese to restructure their own varieties in acquiring SLM; they were,
after all, speakers of larger, socially more prestigious languages in which the SLM
speakers would have been quite competent.'' To return to the original idea by \citet{SmithEtAl2006cll}: Why should Tamil mothers try to make Malay the language of the home? What incentive could they have had? They were speakers of a majority language after all. If, as Smith et al. argue, the Tamil mothers' keeping to Tamil would have posed a problem for the explanation of the maintenance of the language, there is an easier solution than assuming that they chose Malay as a home language: the amount of Tamil mothers was much more limited than what Smith et al. assume. As such, they could not exert as  important an influence. We will see in the following sections that the historical evidence suggests, indeed, that the extent of intermarriage has been overstated in the literature  (Also see Section \ref{sec:slmbg:RelationtotheMoors}).

\paragraph{Sociological evidence}
The historical evidence for intermarriage is said to come from Schweitzer's report and from the Dutch tombos. No evidence is presented as to the minority immigrant language being the superstrate and the local majority language the substrate. No argument is made why children would acquire the mother's non-native variety.

\paragraph{Linguistic evidence}
\citet{SmithRH} notes that Sinhala and Tamil are very similar in their markup, and notes left-branching word order, quotatives, case markers, and similar nominal and verbal categories. Smith notes that Hussainmiya claimed Tamil influence for SLM structures which could also have been of Sinhala influence, and that Jayasuriya claimed Sinhala influence for structures which could as well have been of Tamil influence. As a consequence, he investigates areas where Sinhala and Tamil differ, and finds that SLM is closer to Tamil in the areas he investigated, namely definiteness, number, animacy, and the accusative case. These findings are also used in later publications  \citep{SmithEtAl2004, SmithEtAl2006cll}. This would support a more intensive contact with Tamil speakers.

\citet{SmithEtAl2004} argue that, in the domain of phonology, SLM has a constraint which bans retroflex stops in initial position, just like Tamil, but unlike Sinhala.

\paragraph{Counter-evidence}
The hypothesis hinges on the assumption of heavy intermarriage with the Moor community. As argued above in Section \ref{sec:slmbg:RelationtotheMoors}, the evidence for intermarriage has been grossly overstated. This is not to say that there was no intermarriage. I just fundamentally call into question all the evidence which has been presented for \em heavy \em  intermarriage, which generally does not adhere to the standards of the field. For instance, in \citet[176]{SmithEtAl2006cll} we find again ``intermarriage was first noted by a German observer, Christoph Schwei[t]zer.'' It is true that Schweitzer talked about intermarriage in this passage, but this merits a full quotation to illustrate the quality of sourcing used to defend the Moor intermarriage hypothesis:

\begin{quote}
    Their language is Ambonese, but the majority also speak Malay, Sinhala, Portuguese and Dutch. \el{} The wives, who in part are Ambonese, in part Sinhalese and Malabarian may not say anything [against the stripping of their ornaments]. \citep[106]{Schweitzer1931}, [my translation]\footnote{\el ihre Sprach ist Ambonesisch können aber meistentheils ach Malley- Singules- Portuges- und Holländisch \el die Weiber, welche theils Ambonesin, Singulesin und Malabarin, dörffen nichts darzusagen.}
\end{quote}

There is no explicit mention of Moors in this passage, but Malabarians (a term referring to South Indians, including Tamils) are mentioned, however without mentioning the religion, which could be Hindu or Islam (Moor). As for the languages, we find Malay, Sinhalese, Portuguese and Dutch. Tamil is not mentioned. Now, this is only one passage from a report made by a non-linguist at one point in time.
We cannot exclude that there was intermarriage with Moors, and we cannot exclude that the soldiers' wives spoke Tamil. The source quality is not sufficient for that. But even less can we claim, based on this source, that there \em was \em intermarriage with Moors or that the Malays did speak Tamil. This  is, however, what Smith et al. do.

Another problematic case are the Dutch tombos, which are said to contain evidence for heavy intermarriage.  \citet{Ansaldo2008genesis} checked these records and found only two cases of intermarriage for the Dutch period. This is  surely not enough to support claims of heavy intermarriage. For more discussion about the poor scientific practice with regard to claims about early intermarriage, refer to Section \ref{sec:slmbg:RelationtotheMoors}.

The second assumption is that Tamil was the substrate and Malay the superstrate. This is clearly an attempt to apply the substratist framework to the Lankan setting, and that framework requires that the grammar come from the substrate. Hence, Tamil must be the substrate. This is a linguistic argumentation. If we square this linguistic argumentation with sociology, things get problematic. The substrate is supposed to have low prestige, the superstrate high-prestige \citep[99]{ArendsEtAlEd1995pc}. In the colonial setting, the most likely candidates for superstrate would have been the colonial languages Dutch and English \citep{Ansaldo2005ms,Ansaldo2008genesis}.\footnote{See \citet[25]{Saldin2001} for a similar account.} Next would be languages of law, religion, commerce, and art. Law would probably be  a European language. Religion would be Tamil for the Muslims. Commerce and art would be Tamil or Sinhala. The only sociolinguistically important domain from where Malay could have claimed prestige would be the military. But even there, the colonial languages were present, and used in the higher ranks, so that there is no domain in which Malay would have been the dominant language \citep{Ansaldo2008genesis}.
It is therefore difficult to conceive how the Tamils could have interpreted the Malay language as superior, hence superstratal, and a worthwhile target language to acquire. This applies even more since the Tamil language itself enjoyed high prestige, being the language of the rulers.

\begin{quote}
Kandy was ruled in its last years by a dynasty of Tamil-speaking kings, the Nayakkaras from Madurai in south India.  \el{}
[S]ome of the Kandyan chiefs who were signatories to the Kandyan convention (drawn up when the kingdom fell to the British)
 signed their names in Tamil script rather than Sinhala. \citep[23f]{NissanEtAl1990}
\end{quote}

All this shows that speakers of Tamil had little reason to see their language as devoid of prestige (i.e. substratal) as compared to Malay.\footnote{The prestige situtation may have been different in the genesis of Sri Lanka Portuguese. Portuguese as the language of a colonial power is more likely to have enjoyed high prestige. The Malays however were not in the same position of power as the Portuguese, and it seems clearly inappropriate to equate the sociolinguistic position of the Portuguese with the sociolinguistic position of the Malays.} \citet{Ansaldo2008genesis} furthermore criticizes the use of the term `superstrate' for the language spoken by the immigrants, as this was  a lingua franca with very low prestige to begin with, which does not easily lend itself to superstrate function.

The third assumption is that in bilingual marriages, the children learn the imperfect variety of the husband's language as spoken by the wife, rather than the fluent language of the wife and/or the fluent language of the husband. This would presuppose that the Tamil mothers spoke `broken' Malay to their children, rather than Tamil. The input in `broken Malay' by Moor mothers would have had to be more important than the input in `good Malay' by the husbands and other fluent speakers in the extended family. This is not very convincing as a scenario.


\paragraph{Linguistic counter-evidence}
\citet{SmithRH} identified four areas where SLM was closer to Tamil than to Sinhala: definiteness, number, animacy, and the accusative. Sinhala has obligatory marking of definiteness, while Tamil does not. Smith claims that SLM aligns with Tamil here, but Sections \ref{sec:morph:Indefinitenessclitic} (p. \pageref{sec:morph:Indefinitenessclitic}) and \ref{sec:pred:Nominalpredicate,ascriptive} (p. \pageref{sec:pred:Nominalpredicate,ascriptive})  of this grammar clearly show that indefiniteness marking is obligatory, and thus that SLM is closer to Sinhala here than to Tamil. As for number marking, this is obligatory in Sinhala, and optional on  Tamil inanimates \citep[20]{Lehmann1989}. Smith claims that SLM aligns with Tamil in this domain, but it turns out that number marking is also optional with animates in SLM (Section \ref{sec:func:mod:Plurality}, p. \pageref{sec:func:mod:Plurality}), so that this seems to be a retention of the historic liberal marking of number of Malay varieties rather than copying of the Tamil pattern. As far as the interplay of accusative and animacy is concerned, Smith claims that SLM aligns with Tamil in this domain. In this grammar, it is found that SLM takes a middle position between Sinhala and Tamil as far as accusative marking is concerned. Sinhala has heavy restrictions on accusative marking, Tamil makes more use of the accusative in clearly demarcated areas, and SLM is very permissive. SLM can use the accusative in areas where Sinhala cannot (inanimates), but accusative marking is hardly ever obligatory. This is different from what is found in Tamil, where definite objects are obligatorily marked for accusative. See Section \ref{sec:morph:=yang} (p. \pageref{sec:morph:=yang}) for a discussion of the accusative in SLM.

As far as phonological evidence is concerned, the Tamil constraint against retroflex onsets claimed by Smith is not found as such in SLM. There are both dental onsets (\phontrs{\dentd a:\dentt aN}{come}) and retroflex onsets (\phontrs{\dz a:pur}{oven}) in SLM, which would point  to Sinhala influence in phonology rather than Tamil influence. Actually, voiced dental onsets are a lot less common than voiced retroflex onsets (9 against 42 in Saldin's (2007) dictionary\nocite{Saldin2007Dico}). Furthermore,  Tamil also bans voiced stops from absolute onsets, which is not what we find in SLM. The most salient feature of Tamil phonology, the complete absence of voicing distinctions, has not had any repercussions on SLM phonology, where all traditional voicing contrasts are maintained. If there was indeed heavy Tamil L2-learner influence in SLM, we would expect initial devoicing, devoicing of geminates, and voicing of singleton voiceless stops. None of this is found. If we leave the domain of constraints and turn to segmental phonology, it turns out that the phoneme inventories of SLM and Sinhala are nearly identical, while Tamil is divergent in having no voiced stops, no prenasalized stops, no schwa, but showing retroflex sonorants, which are not found in SLM or Sinhala (Table \ref{tab:Phonemeinventoriescontrasted}).


To sum up, in the domain of morphosyntax we have Sinhala influence in  indefiniteness marking, retention of traditional optional number marking, and a muddled accusative system, which is somewhere between Sinhala and Tamil. In phonology, the phoneme inventories do not suggest Tamil influence, and the purported constraints on the occurrence of initial retroflexes do not hold either. In short, there is no clear sign of exclusive Tamil influence.

\begin{table}
	\subtable[SLM consonant phonemes]{
    \centering
        \begin{tabular}{lccccccc}
                 & labial & dental & alveolar & retroflex    & palatal          & velar & glottal\\
        \hline
        stops    &&&&&&\\
        ~~~voiceless& p   & \dentt{} && \tz        &    c        & k   &   \\
        ~~~voiced   & b   & \dentd{} && \tz        &   \J          & g   &   \\
	~~~prenasalized&\mb& 	     && \ndz        &  \nJ        & \nG &   \\
	nasals      & m   &         &n&             & \ny	  & \ng &   \\
        fricatives  &(f)   &      &        (s)(z)  (\textesh)   && &     & (h) \\
	 approximants & \V  &          &&       &   j         &     &   \\ 
	 liquids & & & r, l &&&&  \\
        \end{tabular}
	}
	\subtable[Sinhala consonant phonemes based on and adapted from \citet{Gair2003, Karunatillake2004}]{
    \centering
        \begin{tabular}{lccccccc}
                 & labial & dental & alveolar & retroflex    & palatal          & velar & glottal\\
        \hline
        stops    &&&&&&\\
        ~~~voiceless& p   & \dentt{} && \tz        &   c         & k   &   \\
        ~~~voiced   & b   & \dentd{} && \dz        &   \J           & g   &   \\
	~~~prenasalized&\mb&\ndentd      && \ndz        &           & \nG &   \\
	nasals      & m   &         &n&               & \ny	  & \ng &   \\
        fricatives  &(f)   &      &        (s)(z)  (\textesh)   && &     & (h) \\
	approximants & \V  &          &&      &   j         &     &   \\
	 liquids & & & r, l &&&&  \\
        \end{tabular}
	}
	\subtable[Sri Lankan Tamil consonant phonemes  based on and adapted from \citet{Zvelebil1966,Suseendirarajah1973phon,Batticaloa,Lehmann1989,Schiffman1999}.]{
    \centering
        \begin{tabular}{lccccccc}
                 & labial & dental & alveolar&  retroflex    & palatal          & velar & glottal\\
        \hline
        stops    		& p   & \dentt{} & \underline{t} & \tz{}        &   c(\J)          & k   &   \\
        nasals      & m   &   &  n      &   \nz           & [\ny]	  &  &   \\
        fricatives  &(f)   &      &        [s](z)  (\textesh)   && &     & (h) \\
	approximants & \V  &  &        &     &   j         &     &   \\
	 liquids &     &               &  r, l & \lz{} &&&  \\
        \end{tabular}
	}
	\caption[Consonant inventories in the adstrates]{Consonant inventories in the adstrates. Parentheses indicate phonemes only found in loanwords. Brackets indicate positional variants which are sometimes included as phonemes.\\ Inventories of SLM and Sinhala are quite similar, while Tamil is divergent. These charts use traditional notation; in the phonology section of this grammar (p. \pageref{sec:phon:Distinctionsinplaceofarticulation}), I propose a reanalysis of the retroflexes.}
  	\label{tab:Phonemeinventoriescontrasted}
\end{table}
 

% Taking a look at historical records of Literary Malay, which show some localisms, we find indications of both Tamil and Sinhala influence. \citet[118]{Hussainmiya1987} notes that in the \em Syair Kisahnya Khabar Orang Wolenter Benggali\em, we find substitution of \graphem{k} for what would be \graphem{g} in Standard Malay. Initial devoicing is a typical feature of Tamil phonology, so that the three occurrences of this localism could be explained by Tamil influence. On the other hand, there are some instances of \graphem{sy} for \graphem{s}. The \graphem{y} could indicate palatalization of /s/. This is even more likely since two of the words stem from Sanskrit, where they had a palatal/retroflex fricative: \graphem{desya} `region' (Skrt. \em de\'sa\em) and \graphem{bahasya} `language' (Skrt. \em bh\=a\dotS\=a\em). These types of fricatives changed to alveolar fricatives in South East Asia, hence the Std. Malay spelling \graphem{bahasa} and \graphem{desa}, respectively. But in Sinhala and Tamil, they merged to the palatal fricative \citep{Gair2003,Karunatillake2004}. So, in Sinhala, we find \phonet{ba:Sa:Va, \dentd e:Saja}, at least among educated speakers.\footnote{Actually, the three sibilants had all merged to /s/ \citep[479]{Geiger1973}, but /\'s/ has evolved into a sociolinguistic marker of register used in loanwords \citep{GairDiglossia}.} It appears that contact with the Sinhalese/Tamil cognates of these words caused the Malays to realize the historical pronunciation of these words, and to emulate it in writing.
% While there are cognates of \em de\'sa \em in Tamil,  \em bh\=a\dotS \=a  \em crucially has no common cognate in Tamil, so that Tamil influence can be ruled out here.


\paragraph{Place in language contact theory}
This hypothesis is well inscribed in substratist theories of creole genesis, although the application of the terms `substrate' and `superstrate'  to the Sri Lankan setting is problematic \citep{Ansaldo2008genesis}.

Furthermore, rather than the standard split grammar/lexicon, SLM seems to have undergone a more fine-grained  mixture of domains \citep{Ansaldo2008genesis}. Parts of verbal morphology are clearly Malay \citep{Slomanson2006cll}, other parts of the grammar are clearly Lankan, and phonology is a mix of Sinhala and Malay features (see relevant chapters in this book).

\subsubsection{The metatypy hypothesis}\label{sec:slmbg:metatypy}
This hypothesis is a variation of the rapid convergence hypothesis discussed above, with the modification that no claims about the speed of development are made, and that Sinhala and Tamil are both accepted as possible sources of grammatical features. \citet{Ansaldo2005ms} contains the first formulation of that hypothesis, while the term `metatypy' is first used in \citet{Bakker2006}. The idea is fleshed out in more detail in \citet{Ansaldo2008genesis,Ansaldo2009book} and \citet{AnsaldoEtAlFCLT}. This hypothesis states that Sinhala and Tamil were adstrates, which both exerted pressure on SLM. The Malays had to be conversant in these two languages of wider communication, which led to the copying of the Lankan structures into Malay grammar.

\paragraph{Argument}
Malay was an esoteric language for the immigrants, in the sense of \citet{Ross1996}, while Sinhala and Tamil were exoteric languages, i.e. languages of wider communication. Sinhala was used to interact with the Sinhalese population, which, at that point in time in the relevant regions would have made up about 80\% of the population. Tamil was used to interact with the rest of the population, but was used more frequently than the pure numbers would suggest because of religious affinities. The grammatical structures of Sinhala and Tamil are very similar.\footnote{Actually, the membership of Sinhala in the Indo-Aryan language family (rather than Dravidian) has occasionally been disputed \citep{Lassen1847,Tennent1859}. See \citet{Geiger1938} for a discussion.} Their typological profiles are close to identical, as noted by \citet{SmithRH}. The Malay language was not a standardized or \em  focused \em variety in the sense of \citet{LePageEtAl1985}. Rather, it  was a \em diffuse \em variety, i.e. a variety where the actual realization of an utterance contained little sociological information as to membership of a linguistic group. The immigrants were all  from different regions, had different mother tongues, and did not attach much importance to the trade language they used to communicate with each other.\footnote{This does not preclude focusing of the language at a later point in time.} As such, this trade language was very variable, and very malleable. This malleable language was then confronted with two adstrates of similar typological markup, which imposed that markup on the immigrants' language. However, while grammar was malleable, this was not the case of the lexicon, which remained by and large Malayic. It can be speculated that the grammar of that language at that point in time was diffuse, with little `acts of identity' \citep{LePageEtAl1985} associated with it, while it was the lexicon that served as an emblem of group membership.

\paragraph{Sociolinguistic setting assumed}
This hypothesis assumes that Malay was an esoteric language and at least one Lankan language was a relevant exoteric language \citep[Tamil in the case of][]{Bakker1995nl,Bakker1996stuf,Bakker2000convergence,Bakker2000rapid,Bakker2006}, although both Sinhala and Tamil would also be possible \citep{Ansaldo2005ms,Ansaldo2008genesis}. It furthermore assumes that the esoteric language was a diffuse variety.

\paragraph{Mechanism}
The mechanism leading to the present state of SLM is argued to be metatypy as described by \citet{Ross1996,Ross1997,Ross2007}.  Bilingual speakers have a tendency to lessen the cognitive load which different grammars impose on them \citep{Nadkarni1975, Ross2007}. In the Lankan context, it seems plausible that the soldiers carried over some Lankan constructions into their Malay speech. Since everybody had some command of the local languages, the speaker could be sure that the hearer would be able to decode the message as intended, even if a Lankan construction was used instead of a Malay one. Furthermore, there was no fear of reprehension for the use of a foreign construction in Malay because of the diffuse status of the code, where no one felt the need to oversee linguistic usage. By and by, the use of Lankan constructions made inroads into the grammar up to the state of the language we find today.

\paragraph{Plausibility}
The Lankan ecology very much resembles other settings where metatypy has been found, so that this hypothesis is plausible.

\paragraph{Sociological evidence}
The esoteric nature of the Malay language can be derived from the confined dwellings in the barracks, with very tightly-knit networks, as described by \citep{Hussainmiya1990}. The exoteric nature of Tamil as language of wider communication is obvious and agreed upon by all authors. The relevance of Sinhala can be derived from the high Sinhalese percentage in the population (80\%). The diffuse nature of the immigrants' language is also uncontroversial \citep[163]{SmithEtAl2006cll}.

\paragraph{Linguistic evidence}
The linguistic evidence for this hypothesis is indirect. SLM exhibits both Tamil and Sinhalese structures, in about the same proportions \citep{AnsaldoEtAl2006SWL,AnsaldoEtAlFCLT}. This suggests that both languages had a comparable status in the genesis of SLM. The most obvious mechanism to accommodate this is to postulate that they were adstrates on a par with each other, which fits the metatypy hypothesis. Other explanations might be possible, but would require a more elaborate apparatus and require more coincidences (e.g. first Tamil influence, then subsequent Sinhala influence of the same importance in the same areas).

\paragraph{Counter-evidence}
Evidence against metatypy could be found in influences which can only be attributed to one language, and not the other. This is the case for lexical influence from Tamil and phonological influence from Sinhala. There are many Tamil loanwords for basic vocabulary terms, like \trs{kattil}{bed}, as well as many animal names (\trs{vanaati}{butterfly}, \trs{vavval}{bat}). On the other hand, there are close to no integrated Sinhala loans. The Tamil loans are phonologically integrated in the SLM system, and also used by speakers who do not know Tamil. The Sinhalese borrowings seem to be nonce borrowings, and the speakers are aware of the code-switch. This difference in lexical influence suggests that at a certain point in time, which was decisive for the vocabulary, Tamil was more important than Sinhala. This point in time was most likely very early, when the Malays had to find new words for unfamiliar concepts (although the concepts `bed',  `butterfly' and `bat' would have existed in South East Asia as well).

Sinhala seems to have had a greater influence on SLM phonology than Tamil, which might have to do with the typologically uncommon nature of Sri Lankan Tamil phonology (three coronal stops, retroflex nasal and lateral, no voicing distinction).
While this shows that Tamil and Sinhala did not have precisely the same function in the genesis of SLM, it does not pose a problem for the metatypy hypothesis because metatypy is only concerned with grammatical structure, and not with the lexicon or phonology \citep{Ross2007}.

\paragraph{Place in language contact theory}
Metatypy is a well-described mechanism of language change in a contact setting. For the Sri Lankan context, it has recently \citep{Ansaldo2009clfet,Ansaldo2009book} been formalized along evolutionary ideas following the ideas of \citet{Croft2000elc} or \citet{Mufwene2001ele}. It has furthermore  been checked against empirical data \citep{AnsaldoEtAlFCLT}. Where the Sri Lankan context differs from more classical cases of metatypy is that the number of languages is greater than 2, a problem shared with the relexification approach (see above). But within the relexification approach, the necessity for two languages is a built-in feature, while for metatypy, it appears to be accidental. The shift in grammatical patterns, which is typical for metatypy, occurs towards the language of wider communication. There is nothing fundamental in the theory of metatypy to preclude change towards two typologically very similar languages of similar use in wider communication at the same time, in this case Sinhala and Tamil. The fact that this theoretical possibility has not been observed earlier is probably due to the very low chance of encountering a trilingual setting where the power relations are `just right' to permit this kind of change.

Metatypy can be seen as the result of speakers adapting their linguistic behaviour in a bilingual setting in order to lessen their cognitive load \citep{Nadkarni1975,Ross2007}. In a 1:1 setting, this reduction of cognitive burden could go either way, but in a 2:1 setting, as in Sri Lanka, there is a clear preference for a `winner-take-all' result, i.e. the development of a 3:0 pattern  as the result of competition and selection among constructions \citep{Ansaldo2005ms}. The social and cognitive mechanisms at work are beyond the scope of this thesis, but are explored in detail in \citet{Ansaldo2009book} and applied in \citet{AnsaldoEtAlFCLT}.
 

 
 
% 
%  \citet{Bakker2000convergence,Bakker2000rapid, BakkerIntertwine}  argues that SLM was a typical creole language
% in the end of the 19$^{th}$ century, isolating, and with mostly SVO
% word order. After the disbandment of the regiment, and with the dispersion of the community to other places than the barracks, Lankan grammar would have made rapid inroads into the grammar of the Sri Lankan Malays.
% 
% 
% 
% 
% While this sounds plausible, one wonders how he comes
% to this conclusion in absence of any data of 19th spoken SLM. 
% Furthermore, VM itself has SVO word order and is close to
% isolating. Bakker affirms that SLM was creolized in the beginning
% of the 19th century. But if VM was already isolating and SVO, the
% fact that SLM was also isolating and SVO in the beginning of the
% 19$^{th}$ century does not allow any conclusion of creolization (If
% ever SLM really was SVO at that time, a fact for that evidence is
% missing.) In the same vein, he argues that SLM has lost person
% inflection at that time. The exact nature of the loss is not
% detailed, which is not astonishing, already VM had no person
% inflection. Bakker draws the conclusion that SLM is a
% ``converted'' language and that it can be shown that the
% conversion took place in one or two generations. Maybe SLM is a
% converted language, but since Bakker doesn't even  attempt to
% adduce any data that would support his claim this question stays
% open.
% 
% \citet{Bakker2000rapid} argues that SLM has undergone radical changes in a few decades in the beginning of the 20$^{th}$ century, comparing the present day language with the language the officers used in their literary writings. But this is comparing a literary variety with a strictly spoken variety. It is not to be assumed that the language the officers used in writing is the same they used when speaking as noted by \citet{SmithEtAl2004}. In South and South East Asia, diglossia is extremely common, and literature will always reflect the high variety. What we know is that the literary documents we have from the end of the 19$^{th}$ century on the one hand and from the beginning of the 20$^{th}$ century on the other hand, differ considerably, the latter showing influence from lankanized SLM \citep[14]{SmithEtAl2004} while the former do not. Smith claims that this shows that 19$^{th}$ century colloquial SLM was already very different from the high variety used in literature. Otherwise the change would not have been so drastic. Bakker on the other hand argues that SLM ``converted'' within one or two generations, which would also fit well with the rapid change shown in the literary data.
% 
% But on the basis of literary language alone, it is difficult to tell which hypothesis is right. We do know for sure that the literary variety from the 20$^{th}$ century is more removed from standard Malay than the literary variety from the 19$^{th}$ century. But we have no hints on whether each spoken variety was equidistant to its literary counterpart, or whether the distance between the two diminished over time. One can of course speculate that the literary variety reflects the colloquial variety, which would strengthen Bakker's case, or that the decline of the literary variety could only take place because spoken SLM was already very much lankanized before, which would be Smith's line of argumentation. But in absence of more concluding data, the decision cannot be made.
% 
% \begin{quotation}
%     ``If [the accusative marker] does indeed derive from a feature of Javanese or Betawi \el, there is all the more reason to expect an early appearance in SLM, as Java-based retional peculiarities could not have lasted long in Sri Lanka, particularly once the British took over and stopped recruitment from Java.''\citet[19]{SmithRH}
% \end{quotation}

%
%
%\begin{quotation}
%   ``What  kind of Malay was spoken by the Javanese and other diverse peoples who were brought to Sri Lanka from the Dutch East Indies? It could hardly be anything other than the Vehicular Malay which had for centuries served as a lingua franca of [the SEAA]. This lingua franca was/is based in spoken low Malay rather than the high varieties which gave rise to the national standards \el. The term ``Bazaar Malay'', which has also been used to refer to this variety, implies a much narrower domain of use than the languages actually had (and continues to  have) and is also used in other senses (Paauw 2003); \nocite{Paauw2003} hence I believe the term Vehicular Malay is to be preferred. Vehicular Malay was not and is not a monolithic variety, and there was clearly some regional differentiation. In particular, among the varieties of Vehicular Malay brought to Sri Lanka, Betawi, the low Malay of the Batavia area, would have been strongly represented, since Batavia served as a staging ground for troops being sent to Ceylon.  ''\citet[14f.]{SmithRH}
%\end{quotation}

%\begin{quote}
%   ``It appears, however, that the early Indonesian migrants, drawn from such varied eastern races had shed their different identities even before they were introduced to Sri Lanka, and had evolved intro a single identity through the use of the Malay language, which served in uniting all these different national groups''\citet[11]{Hussainmiya1987}
%\end{quote}

% 
% \subsection{Dating of Influences and Influencing from Input Languages} \label{sec:slmbg:DatingofInfluencesandInfluencingfromLankanLanguages}
% It is sure that the language spoken by the Malays in the 17$^{th}$ century has changed considerably in the last three centuries to become the language that is spoken today in the SLM communities. The language has become considerably less Malayo-Polynesian and considerably more Lankan. Not much research has been undertaken on  the timing of these changes.
% 
% To get from the initial state of the language (S_0) to  the state of today S_{now}\footnote{Members of a community do not all make use of recent changes. Some are    quicker to adopt them, some are more conservative %\citet{Milroy}.
% Thus, S_0 as well as S_{today} are generalizations.}
% ,there are a basically three possibilities:
% \begin{enumerate}
%         \item the language changed steadily over time
%         \item the language changed abruptly at one point (or several) in time
%     \item there were times of quicker language change and times of less change
% \end{enumerate}
% 
% \begin{figure}
%     \centering
%         \includegraphics{pics/lgchggrph}
%     \caption{Three types of language change}
%     \label{fig:lgchangegraph}
% \end{figure}
% 
% Given the vicissitudes of the Malay history in Ceylon, a steady change does not seem probable.
% If there was a very abrupt language change, it must have been motivated by sociohistoric events of magnitude. One of these events is certainly the arrival of the first Malays in Ceylon, other possibilities are the eviction of the Dutch by the British, the disbandment of the Regiment and independence of Ceylon.
% 
% How likely was abrupt language change at these four points in time? The first settlers to arrive were exiles. Exiles do not leave their country of origin willingly. They prefer to stay in their homeland. The exiles (\em bannelingen\em) in Ceylon were insurgent leaders from the SEAA. They fought against Dutch rulership and for independence. As such, they are likely to value their own culture very highly. Furthermore, they hope to set again foot in their homelands. This makes them not very likely to innovate or adopt linguistic changes innovated by others.
% 
% The exiles were the Malay `elite' in Ceylon.  They numbered about
% 10\% of the population. At least they would not have undergone a
% major and abrupt linguistic change. They were serving as a
% guideline for the other groups, with whom they interacted. Some
% change is likely to have occurred, but not of a very abrupt
% nature.
% 
% The British takeover from the Dutch was rather smooth for the Malays, who quickly reassumed their old role as soldiers. The language of the military changed from Dutch to English. The two European languages being fairly similar in structure, and fairly remote from Malay or either Lankan language, a major impact of that is not to be expected.
% 
% The fading out and ultimate disbandment of the Regiment caused major distortions in Malay
% society. They lost their hearth in the Regiment, their social
% cohesion and many moved from  the garrison town to other places
% to find work. This changed the sociolinguistic setting in a more
% catastrophic way then the British takeover. With the
% sociolinguistic setting totally turned upside down, major repercussions on the language cannot be excluded.
% 
% \begin{quote}
% 	``networks chiefly constituted of strong ties support minority languages resisting institutional pressures to language shift; but when these networks weaken, language shift is likely to take place''\citep[124]{MilroyEtAl2003}
% \end{quote}
% 
% 
% Independence from the United Kingdom, nationalist movements and the promotion of the infamous ``Sinhala only'' law in 1956   \citep[35]{NissanEtAl1990} also had an impact on the sociolinguistic setting. School lessons were now in Sinhala, so the Malay children had to learn this language. They could no longer attend the English medium schools. On the other hand, English was still the language of economic ascendance. Many Malay families opted thus for English as the language spoken at the home to provide their children with good future opportunities \citep{Saldin2003}. This was to the disadvantage of Malay, and the source of language attrition that we observe today. Linguistic change in the post-colonial era was surely major, but not abrupt as it was only the younger generation that lost knowledge of Malay, while the parents still were in command of the language, but limited its use. Major Sinhala influence is likely to date from that period.
% 
% I conclude that no abrupt change happened at any of these four
% occasions (first settlement, British rule, disbandment of the
% regiment, independence.)
% 
% By exclusion of the first two possibilities, the change must have been gradual, with periods of heavy change alternating with periods of less change. The periods in which major changes are likely are the following:
%  
% \begin{enumerate}
%     \item the late 17$^{th}$ century: consolidation of a Lankan vehicular Malay
%     \item 1873+: dispersal of the community, loss of cohesion: major change
%     \item 1950+: forced introduction of Sinhala: major change
% \end{enumerate}
% 
% This view is supported in \citet{Ansaldo2005ms}.

 

\subsection{The linguistic ecology}
In Section \ref{sec:slmbg:SourcesintheSEAA} we have seen the diverse influences from the South East Asian varieties in different periods of history. In this section, I will briefly sketch the present linguistic ecology of Sri Lanka to shed light on which languages the Sri Lankan Malays encounter in their daily life.

The majority language of the island is Sinhala. Sinhala is a diglossic language \citep{Gair1968,
Gair1985diglossia,
Paolillo1992,
Paolillo1997},
with important differences between the high and the low variety. These differences cross-cut the spoken/written distinction, i.e. it is possible to combine high or low with spoken or written expression \citep{DeSilva1974,Gair1985diglossia,Paolillo1997}. The high variety is described by
\citet{Gunasekara1891},
\citet{Geiger1938,
Geiger1973},
and \citet{GairEtAl1974literary},
while the low variety is described by
\citet{Chater1815,
FairbanksEtAl1968,
Matzel1983,
GairEtAl1997,
Dissanayaka2003,
Jayawardena2004}
and \citet{Karunatillake2004}.
Linguistic analyses of various topics of Colloquial Sinhala are compiled in \citet{Gair1998} and \citet{Henadeerage2002}. The low variety is further divided between `Upcountry' and `Lowcountry' dialects, but the two are mutually comprehensible. The further dialectal subdivision is unclear and underresearched \citep[xii]{Jayawardena1996}. The descriptive material available for Colloquial Sinhala is  good and covers the speech of the groups the Malays have contact with.

The next biggest language of Sri Lanka is Tamil. While there is abundant literature on Tamil, it very often does not address the varieties relevant in the Sri Lankan context. Like Sinhala, Tamil is a diglossic language, with Literary Tamil being distinguished from diverse forms of Colloquial Tamil \citep{AnnamalaiEtAl1998}.  Colloquial Tamil can be divided into Indian Tamil and Sri Lankan Tamil. Spoken Indian Tamil is well described
\citep{Pope1867,
Arden1934,
Beythan1943,
Asher1982,
Lehmann1989,
Schiffman1999,
AsherEtAl2002,
AnnamalaiEtAl1998},
but differs from Sri Lankan varieties. Indian Tamil is present in Sri Lanka via media (television, films), and via the Estate Tamils, who came during the British period to work on the tea plantations. `Media Tamil' is quite homogeneous, but Estate Tamil is likely to still reflect the dialectal differences of the Indian immigrants from different regions of Southern India. The Malays do watch Tamil movies, and were in contact with Estate Tamils when they were overseers. This means that these two varieties of Indian Tamil can be considered as contact languages.

Dialectal variation in Tamil is underresearched \citep{Pillai1986}, especially in what concerns Sri Lanka.
\citet{Zvelebil1960,
Zvelebil1966}
distinguishes Jaffna Tamil, Trincomalee Tamil, Batticaloa Tamil and General Ceylon Tamil  on the basis of the phonology and morphology of a handful of speakers. \citet{Suseendirarajah1967,Suseendirarajah1973phon} adds some more description of the Tamil of Jaffna. The dialects of Jaffna and Trincomalee are not important for the analysis of SLM since the Malays had little contact with those towns.\footnote{There are some Malays in Trincomalee, but it is unclear whether they still speak Malay \citep[426]{Hussein2007}.} `General Ceylon Tamil' is more promising, as is Batticaloa Tamil \citep{Zvelebil1966, Suseendirarajah1973pron,Batticaloa}, with an important percentage of Muslim speakers. With the exception of Batticaloa, all these descriptions cover Hindu/Christian varieties. As for dialects on the other side of the island, \citet{Bonta2004,Bonta2008} argues that varieties of Negombo Tamil have converged towards Sinhala and have lost verb agreement.

The Moors have to be divided on sociological grounds into `Coast Moors' and `Ceylon Moors', of which the Ceylon Moors speak a Lankan variety while the Coast Moors, hailing from India, speak an Indian form of Tamil \citep[25]{Nuhman2007}, also discussed in \citet{Hussein2007}.
\citet[71ff]{Nuhman2007} distinguishes Sri Lanka Muslim Tamil from the Hindu/Christian varieties. He further divides North-Eastern Muslim Tamil from South-Western Muslim Tamil \citep[also see][45]{Hussein2007} and argues that verb agreement was lost in the South-Western variety as well. He furthermore gives an overview of some additional phonological, morphological and lexical characteristics of Muslim Tamil. To sum up, overall description of Tamil is very good, description of Sri Lankan Hindu/Christian Tamil is fair, and description of Muslim varieties is seriously wanting, as is a description of Estate Tamil varieties. Unfortunately, it is the Muslim and the Estate varieties which are especially interesting as contact languages for the Malays.

Another important language is, of course, English, in its Sri Lankan shape \citep{Kandiah1996}, while the  oldest inhabitants of Sri Lanka, the Veddah, and their language \citep{DeSilva1972}, have not been in extensive contact with the Malays.

\subsection{Heydays}
Different languages had different impact in different periods on the formation of SLM, as is common in contact language formation \citep{Arends1989,Roberts2004phd,Lim2007merger}. To go by chronological order, Portuguese was the first language to have the possibility to make an impact, although there is little evidence for that. The only item which could point to a Portuguese influence specific to Sri Lanka is \trs{salba}{escape}, which does not seem to exist in the SEAA. There are numerous other Portuguese words in SLM, but these are either also found in the SEAA (e.g. \trs{greeja/gereja}{church}) or in Sri Lanka (\trs{kusini}{kitchen}, Sinh. \em kussiya, \em Tam. \em kusini\em), so that it cannot be established that they were adopted in Sri Lanka during Portuguese rule. They might have entered the language later, through new immigrants bringing words like \em greeja \em from the SEAA, or through adstrate influence bringing words like \em kusini\em. After the Portuguese were ousted, their influence ceased. Next came the Dutch, who also failed to have a huge impact on the language. No Dutch word seems to be exclusive to SLM; they are all shared with other varieties of Malay and the local languages (like \trs{kaamar}{room}, Std. Malay \em kamar\em, Sinh. \em kaamaraya\em). During Dutch rule, many immigrants came from Java and other islands of the SEAA.
There is a strong trace of them in SLM, with 11,7\% of the vocabulary being traceable to them exclusively \citep{Paauw2004}. The influx from island immigrants ceased in 1819, and the immigration from peninsular Malays grew stronger. This new wave of immigrants, however, did not have an important impact. Only 0,6\% of SLM vocabulary can be traced to them \citep{Paauw2004}.
During the period of the regiment, Literary Malay was taught in the regiment schools, so  this language could have exerted influence on the spoken variety, although it is unclear whether this indeed materialized. When the regiment schools were closed, Malays switched to English education, and would cling to this medium of instruction until independence. English was probably already a major factor, but schooling in English increased the impact of this language from the late 19$^{th}$ century onwards. After Sinhala became the medium of instruction in government schools following independence, Sinhala increased its influence on SLM, and continues to do so up to today, due to the economic necessity to command this language.

As for Tamil, it is likely that it had an influence in the early days because many basic terms like \trs{kattil}{bed} are borrowed from this language, as are many names for animals. It is likely that the variety of Tamil was Muslim Tamil due to the shared religion; there is close to no evidence for prolonged contact with Hindu or Christian Tamils. The relative strength of Muslim Tamil influence over time is difficult to assess. In the beginning, when there were no Malay mosques, it was probably stronger, but as soon as Malay mosques were established and Malay life evolved around the military cantonments, Moor influence is likely to have decreased. The growing association with the Burghers and the adoption of a more intellectual and Westernized lifestyle in the 19$^{th}$ century suggest that the cultural ties between Malays and Moors became weaker in that period than what they had been before. But at the same time, the beginning of the estate economy, which employed many Malays as overseers over South Indian plantation workers, meant that the Malays had once again a good motivation to be conversant in Tamil. While the prestige of Estate Tamil was of course not very high, the relative size of labourers against overseers might have made up for that. As of today, the influence of Estate Tamil on Sri Lanka Malay (at least the Upcountry variety) remains understudied. With the orientation of Malays towards service economy, the contact with Estate Tamil decreased.

To sum up, most linguistic influence came, peaked and possibly disappeared. This is the case for the colonial languages as well as the languages of the immigrants, Estate Tamil and  Sinhala (which came to stay). An exception to this `peak' pattern is Muslim Tamil, which was probably present during the whole period in unclear proportions, possibly with a dip during the period of the regiment.

\subsection{Envoi: Against reductionism}
From the above descriptions, it should be clear that the linguistic history and ecology of Sri Lanka Malay are very complex. Sri Lanka Malay is definitely not a simple combination of Malay and Tamil and/or Sinhala. While this characterization could be used as a shorthand, it oversimplifies the diachronic complexities. To start with the immigrants, they surely spoke more than one language, and `Malay' is probably a misnomer to designate most of their varieties, if used without a qualification (Moluccan Malay, Jakartan Malay, Tangsi Malay etc.). It is not clear whether we can speak of `the immigrants' language' at all, given that the soldiers hailed from very diverse parts of the South East Asian Archipelago. Maybe there was a language which all of them shared, but it is more probable that that was not the case.\footnote{It should be borne in mind that Bahasa Indonesia, the national language of Indonesia, is a 20$^{th}$ century creation and is not spoken natively by a majority of the population \citep{Steinhauer2005}.} It is likely that different kinds of accommodations  were used on the part of a speaker to facilitate communication with addressees from different ethnolinguistic backgrounds, and that the `Malay lingua franca' was perpetually reinvented, very much like its Mediterranean counterpart \citep{Selbach2007,Selbachinpreparation}.

In Sri Lanka, Sinhala can be seen as a relatively homogeneous language. This is not the case for Tamil, where Hindu/Christian varieties and Muslim varieties differ considerably \citep{Nuhman2007}, and the especially interesting Muslim varieties show diatopic variation. Estate Tamil also differs from estate to estate depending on the origin of the workers. Just as we must ask `Which Malay?' when talking about the immigrants' language, we must ask `Which Tamil?' and `Which Sinhala?' when speaking about the language of the people they met in Sri Lanka.  Reconstruction of historical contact settings is notoriously difficult \citep{RenfrewEtAl2000,Lim2007merger,Ansaldo2009clfet}, but it is probable that the answers to these questions will not be `that Tamil' but rather `these Tamil varieties, those Malay varieties, those other Sinhala varieties'. Only when we have a full and precise picture of the input varieties, and the end product, will we be able to make solid analyses of the evolution of Sri Lanka Malay.

Diachronic language contact studies investigate the development of a new variety out of the contact between two or more prior varieties. The more information we have on the historical languages and the present day language, the more precisely we can trace the development of the language. As for Sri Lanka Malay, we have little information about the historical varieties, and not too much information about the current variety. While the Sinhala part seems OK, we need a better description of historical Trade Malay varieties,\footnote{After writing this chapter, I received a copy of \citet{Paauw2008phd}, which is a big step forward into that direction.} historical Muslim Tamil varieties, and present day Sri Lanka Malay to make solid analyses of the development of Sri Lanka Malay.
I will not be able to shed much light on the historical varieties in this thesis, but in the following chapters, I hope to advance our knowledge of the current shape of Sri Lanka Malay and  add the contemporary piece to the diachronic jigsaw puzzle.

